\documentclass[ecta,nameyear,draft]{econsocart}
\usepackage{xr-hyper}
\usepackage{amsmath,amsfonts,amssymb,amsthm,graphicx,booktabs,array,longtable}
\usepackage{algorithm,algpseudocode,bm,enumitem}
\usepackage[T1]{fontenc}
\usepackage{mathpazo}
\usepackage{microtype}
\usepackage{needspace}
\RequirePackage[colorlinks,citecolor=blue,linkcolor=blue,urlcolor=blue]{hyperref}
\startlocaldefs
\newtheorem{theorem}{Theorem}
\newtheorem{proposition}[theorem]{Proposition}
\newtheorem{lemma}[theorem]{Lemma}
\newtheorem{corollary}[theorem]{Corollary}
\newtheorem{assumption}{Assumption}
\theoremstyle{definition}
\newtheorem{definition}{Definition}
\theoremstyle{remark}
\newtheorem{remark}{Remark}
\newtheorem{example}{Example}
\newcommand{\E}{\mathbb E}
\newcommand{\R}{\mathbb R}
\newcommand{\HH}{\mathcal H}
\newcommand{\LL}{\mathcal L}
\newcommand{\Act}{\mathcal A}
\newcommand{\sg}{\operatorname{sg}}
\newcommand{\tr}{\operatorname{tr}}
\newcommand{\Lip}{\operatorname{Lip}}
\DeclareMathOperator*{\argmax}{arg\,max}
\DeclareMathOperator*{\argmin}{arg\,min}
\endlocaldefs

\let\NBOLongtable\longtable
\def\longtable{\Needspace{7\baselineskip}\NBOLongtable}

\setlength{\emergencystretch}{2em}
\externaldocument{revisions/2026-10-07-r41/build/ECTA_from_historical_article}[historical_article.pdf]
\externaldocument{revisions/2026-10-07-r41/build/ECTA_from_historical_supplement}[historical_supplement.pdf]
\externaldocument{revisions/2026-10-07-r41/build/ECTA_from_applications}[applications.pdf]
\externaldocument{revisions/2026-10-07-r41/build/ECTA_from_supp}[supp.pdf]
\externaldocument{revisions/2026-10-07-r41/build/ECTA_from_response}[response.pdf]

\makeatletter
\def\copyright@text{Prepared for referee review}
\makeatother
\makeatletter\def\copyright@text{Prepared for referee review}\makeatother
\begin{document}
\begin{frontmatter}
\title{Neural Bellman Operators}
\runtitle{Neural Bellman Operators}
\begin{aug}
\author[id=au1,addressref={add1}]{\fnms{Qian}~\snm{QI}\ead[label=e1]{qiqian@pku.edu.cn}}
\address[id=add1]{Peking University}
\end{aug}
\begin{abstract}
This paper develops Neural Bellman Operators for policy evaluation and feasible improvement in controlled economies. Centered continuation errors connect learned values to full-policy loss. A direct nonlinear certificate evaluates each trained network against the Bellman operator containing its own future network, including all-state residuals, selected actions, and terminal error. A separate execution account covers rounded state acquisition, transitions, and costs. A quadratic construction connects hidden-factor updates to economic tolerances; a pressure experiment activates native precision escalation and retains rejected attempts. In nonlinear investment economies, direct neural policy-loss bounds at declared states are smaller than verified gains from reoptimization after an investment-price change. Contemporaneously optimized spline policies remain competitive and often have strictly smaller cost. A dimension-explicit extension covers continuous innovations and accounts for verification and policy-query work. The controlled-economy framework and complete applications are retained with their application-specific hypotheses.
\end{abstract}
\begin{keyword}
\kwd{Dynamic programming}\kwd{Neural approximation}\kwd{Recursive utility}
\kwd{Endogenous preferences}\kwd{Viscosity solutions}\kwd{Dynamic games}
\end{keyword}
\end{frontmatter}
\section{Introduction}\label{sec:r37intro}
Dynamic economic decisions depend on continuation values that are costly to
recompute and change when future policies, technologies, or preferences
change. A neural approximation can make repeated evaluation inexpensive, but
neither a small fitting loss nor a successful current decision establishes
that the resulting dynamic policy is accurate. The relevant questions are
which continuation errors affect decisions, how those errors propagate
through the policy's own future, and what work is needed to construct and
verify a neural continuation at the required economic accuracy.

This paper develops Neural Bellman Operators to answer these questions within
a single policy-evaluation and feasible-improvement framework. The original
object is a controlled economy with a policy-specific continuation, including
controlled diffusions and recursive preferences. We retain the distinction
between evaluating a fixed policy and optimizing against that evaluation.
Current actions are varied over their actual feasible set; the future policy
is held fixed while its continuation is evaluated. Boundary, utility-domain,
and comparison conditions belong to the economic problem and are not
replaced by a neural parameterization. Two distinct realizations now give constructive policy certificates: a recursive capital economy with Gaussian innovations and vector controls, and a nonlinear compact-state investment economy with constrained actions. The latter admits genuinely nonquadratic continuations and trained ReLU policies. Both are applications of the same operator framework.

\paragraph{The current operational chain.}
Section~\ref{sec:r41direct} makes the stored trained network itself the
Bellman object and distinguishes original-model policy loss from numerical
execution cost. Section~\ref{sec:r41evidence} reports its direct certificates,
contemporaneously reoptimized comparators, and precision-pressure paths.
The dimension and continuous-innovation extension in
Theorem~\ref{thm:r41dimension} exposes verification cost explicitly.
The earlier comparison and training results below remain in force under
their stated assumptions. The scalar nonlinear execution does not stand in
for a coupled high-dimensional performance experiment.

\paragraph{Contribution and roadmap.}
Section~\ref{sec:r39nonlinear} proves an all-adapted-policy comparison and a finite nonlinear construction with explicit off-grid and continuous-action allowances. It also supplies a breakpoint verifier for trained hidden units. Section~\ref{sec:r39floating} certifies the entire native matrix step and accounts for evaluation and verification work. Section~\ref{sec:r39evidence} uses the complete factorial design and develops a nonlinear investment-price counterfactual. References below to the eight earlier exact services describe a historical construction audit; their joint adaptive-cached contrast does not identify a precision effect. The new factorial does not establish a stable adaptive timing advantage.


The first contribution identifies the decision-relevant error. An additive
level shift in a continuation does not change an optimizer. The appropriate
local error is therefore an oscillation across feasible actions, together
with any optimization error and any discrepancy between the evaluated and
implemented futures. This distinction matters for both theory and numerical
diagnostics: a method may have a large level error but a small decision loss,
or a small average fitting loss and an inaccurate action. The centered bounds
also explain why a certificate for a current action cannot simply be reused
as a certificate for a changed future.

The second contribution gives a constructive route from trained hidden
weights to full-policy accuracy. In the recursive capital economy, an
independent coefficient comparison applies on all real states and against
all real vector controls. The comparator is every adapted policy with finite
recursive cost, not merely policies in the fitted class. Positive Gaussian
moment margins control the propagation of quadratic and state-independent
errors separately. A primitive envelope allocates policy error before the
neural targets are supplied. Cold construction, a gradient--curvature
criterion, and basin-certified warm construction then provide explicit routes
to the required continuation accuracy. These statements concern an explicitly
trainable square-activation continuation; every entry of its hidden factor
may change. They are not fixed-feature regression guarantees applied to an
unrelated jointly trained network.

The third contribution joins that construction to finite-precision reuse.
A changed continuation may invalidate the old action while leaving the old
factor in an admissible training basin. We derive a two-sided error budget
that pays separately for factor fitting, own-policy evaluation, and the actor
equation. It also checks the magnitude of the fitted coefficient, preventing
an overshooting factor from invalidating the primitive envelope used to prove
the policy bound. Within this budget, a precision-adaptive refresh rule permits
rectangular factors and noncommuting starts. Its relative Gram enclosure
contracts according to $\rho_{j+1}=\rho_j^2$ when the entire implemented
step error receives the stated allowance. A fixed target inverse is constructed
at most once per fit; subsequent applications and checks remain charged.
The resulting theorem specifies both a finite update cap and a sum of stored
fractional precisions. It does not presume that ordinary floating-point
matrix operations have zero error.

These results supply a positive neural construction theorem, not a claim
that neural methods must dominate every structural solution. We distinguish
three statements throughout the paper. A policy certificate establishes the
accuracy of a returned policy. A training theorem establishes that the declared
neural construction can produce an admissible continuation under stated
primitive and initialization conditions. A complete-work comparison establishes
which executed procedure was less costly at a common economic target. None
of these statements alone establishes the other two. In particular, a Gaussian
quadratic economy has a strong structure-exploiting competitor, and we retain
it rather than make the comparison artificially favorable to learning.

The numerical evidence makes these distinctions observable. The earlier
finite-law study establishes valid current-decision certificates and detects
continuation changes, but its adaptive neural rule is more costly than fresh
neural refitting in all eight reported dimension--volume cells; cached
simulation is the least-cost fully certified procedure in those cells. Those
outcomes and the complete original economic applications remain part of the
publication. Subsequent continuous-support, vector-control studies execute
full-policy comparisons, including a 72-service recursive-risk catalogue.
The new frozen comparison isolates construction, precision, and target reuse
on eight exactly specified services. All eight satisfy the same full-policy
tolerance. The earlier joint adaptive-cached implementation reduces stored precision and has smaller recorded complete-service clocks than the two inherited neural constructions in both specified regimes; that contrast alone does not identify a precision effect. The structural procedure is faster still.
These are deterministic construction comparisons, not estimates of a
population probability of neural superiority.

The economic interpretation likewise depends on the target. A current
withdrawal response conditional on a specified future policy differs from
a counterfactual in which all future controls are reoptimized. We retain
the former finite-law comparisons with their original interpretation and
use the returned policy's own continuation in every date of the latter
recursive construction. The capital primitives are specified theoretical
economies rather than empirical estimates. The separate applications to
endogenous preferences, recursive utility, temporal selves, and games retain
their own utility domains and deviation or equilibrium conditions. The
capital certificate does not silently discharge those different obligations.

The article proceeds from the economic model and Neural Bellman Operator
to decision errors, full-policy comparison, constructive training, and
precision-adaptive reuse. The technical supplement supplies the proofs and
supporting training criteria. The complete previous article, its supplement,
and all economic applications are preserved as identified companions, with
cross-document references. This organization permits the current argument
to be read in economic and mathematical order without discarding earlier
results or requiring the reader to reconstruct the development from branch
history.

\subsection{Relation to the literature}
Neural Bellman Operators build on numerical dynamic programming and function approximation. Projection and related methods exploit economic structure to approximate a Bellman equation; see \citet{judd1998}. Smooth neural networks can approximate derivatives under representation hypotheses \citep{hornik1990}. These results distinguish representational capacity from the accuracy achieved by a finite architecture and a particular optimizer.

The differential critic is related to the neural PDE approach in \citet{sirignano2018}. Their DGM formulation includes differential and boundary conditions and treats HJB equations. The deep BSDE method of \citet{han2018} also includes an HJB application. Our object combines a specified policy continuation, a feasible actor, and an economic error account, and asks when the continuation representation is useful across decision queries. An analytical maximizing action is useful when available; it is not a universal distinction between neural PDE methods. Recursive dependence on the value is a nonlinear zeroth-order term, whereas full nonlinearity of a PDE concerns its highest derivatives.

Recursive utility separates intertemporal substitution from risk aversion \citep{epsteinzin1989,duffieepstein1992}. Utility-process selection and transversality matter in its infinite-horizon formulation \citep{herdegen2021}. Temporal selves instead require an equilibrium formulation \citep{bjork2016}. The stochastic approximation literature provides convergence results under explicit step-size, stability, and limiting-dynamics conditions \citep{borkar2024}. We retain those distinctions when describing a finite neural computation.

Section~\ref{sec:model} defines the economic object and Section~\ref{sec:r10continuous} specializes it to the capital economy. Section~\ref{sec:r15mechanism} connects continuation error with economic gain. The computational design in Section~\ref{sec:r16design} distinguishes continuation reuse, conditional payoff replication, and mechanism assessment. Section~\ref{sec:economic-scope} states the other economic applications. General verification results and their full proofs remain in the technical supplement, beginning with Section~\ref{sec:theory}; the application companion is a current part of the paper.

\subsection{Continuation approximation and the contribution of NBO}
A reusable continuation is not a new principle of dynamic programming.
Postdecision-state approximation separates a controlled transition from the
subsequent uncertainty and is a standard device in approximate dynamic
programming \citep{powell2011}. Regression can estimate a continuation relevant
to an economic stopping decision, as in \citet{longstaff2001}. Batch fitted
value methods turn transition data into successive supervised regression
problems \citep{ernst2005}; approximate policy iteration likewise separates
an evaluation error from an imperfect improvement step
\citep{scherrer2015}. These methods establish the appropriate comparison
class for NBO. The relevant question is not whether continuation fitting is
possible, but which fitted object, decision rule, and verification account
are useful for the stated economic exercise.

The actor--critic distinction is also established. In
\citet{konda1999}, a critic supplies information for a parameterized policy
update under explicit stochastic-approximation conditions. NBO uses a scalar
policy continuation and its derivative channels in a feasible economic
improvement map. The scalar representation couples value and gradient
predictions through one potential; a separately fitted vector field need
not have this property. That structural distinction is testable but does
not imply that a scalar neural network must outperform a costate surrogate,
a direct actor, or a conventional scalar regression. The menu's vector
baseline and the new quadratic, radial-basis, and shared-actor baselines
therefore remain in the primary comparisons.

The reuse comparison is closely related to simulation optimization. A cached
sample-average objective can be optimized at many current tasks. Reusing
innovations is legitimate; reusing their payoff at a different postdecision
state without reevaluation is not. The new experiment gives cached simulation
and regression the same finite-law information. The regressions share their
value and gradient labels, and the shared actor receives the same task
variables. Full-support enumeration is a particularly strong comparator
when the law has only sixteen atoms. Its tractability is disclosed rather
than suppressed to manufacture an advantage for approximation.

Several mathematical components are standard. Centering removes constants
from action comparisons; the squared-error difference cancels common noise;
shortest-path closure solves difference constraints; and a strong-concavity
bound converts a first-order residual into an objective gap. For stochastic
stopping, time-uniform confidence sequences provide a general route to valid
repeated assessment \citep{howard2021}. We do not claim any of these principles
as a new theorem of dynamic programming, graph theory, or probability. The
contribution developed here is their explicit economic linkage: which
continuation is reusable, how approximation and changes in future primitives
enter decision loss, which computed action is certified, and which work is
actually paid before that certificate is returned. The deterministic finite-law
experiment uses no confidence sequence and spends no additional probability.
Its scalar-interval guarantee and the earlier statistical menu guarantee
are different applications of that linkage.

% R15 component; only input paths and enumerated location/grammar prose are edited.
% Source: revisions/2026-10-04-r15/manuscript/model.tex; commit 1cb3cc9efe135966ec228dfaf51c6841a6dfc97c.
% Generated by code/integrate.py; edit extraction rules there.
% Reviewed source: ECTA.tex, line 75, commit f5021cefa71492babfcbfa580e0c984f59a9de26
\section{The Controlled Economy}\label{sec:model}
Let $W$ be an $m$-dimensional Brownian motion on a filtered probability space satisfying the usual conditions. The state $S=(u,X)$ contains internal and external state variables. An action $a=(\theta,c,\pi,\ldots)$ combines internal adjustment and external choices. The controlled state evolves as
\begin{equation}\label{eq:state}
 dS_s=b(s,S_s,a_s)\,ds+\sigma(s,S_s,a_s)\,dW_s,
 \qquad a_s\in\Act(s,S_s).
\end{equation}
The compact feasible correspondence $\Act$ is part of the model. Actions are progressively measurable and must generate a well-defined state process satisfying the declared state restrictions. State viability does not follow from bounded actor outputs alone.

We work initially on a finite horizon $T$. A computational domain $D$ is either invariant, equipped with a specified reflection mechanism, or stopped on exit. In the stopped case let $\tau=T\wedge\inf\{s\geq t:S_s\notin D\}$. The payoff $g(\tau,S_\tau)$ is specified on the entire parabolic stopping boundary, including the terminal face. For a fixed admissible policy $a$, utility is the policy-specific solution
\begin{equation}\label{eq:bsde}
 Y_s^a=g(\tau,S_\tau)+\int_s^\tau f(v,S_v,a_v,Y_v^a)\,dv
                    -\int_s^\tau Z_v^a\,dW_v.
\end{equation}
Adjustment costs are included in $f$. For example, time-additive utility has $f(s,S,a,y)=r_0(s,S,a)-\rho y$, where $r_0$ includes all current resource and adjustment costs. Define $V^a(t,x)=Y_t^{a,t,x}$ and $V^*(t,x)=\sup_a V^a(t,x)$. The value inside a candidate policy's recursive objective is its own utility process, not the unknown optimum inserted into every candidate objective.

\begin{assumption}[Finite-horizon comparison class]\label{ass:model}
The controlled state problem is well posed for the admissible policies used below. The coefficients are continuous, locally Lipschitz in the state uniformly over actions, and bounded on the compact stopped domain. The payoff is continuous and bounded. On an invariant utility interval containing all policy values, $f$ is continuous in its arguments and uniformly Lipschitz in utility with constant $L\geq0$. The admissible class is stable under concatenation and has the measurable-selection properties needed for dynamic programming. The resulting HJB problem and the fixed-policy problems obey comparison in the bounded continuous class with the specified boundary conditions.
\end{assumption}

The comparison requirement identifies both a function class and a boundary problem. It is not asserted merely because the primitive coefficients are continuous. On a bounded, regular Dirichlet domain, a boundary barrier is one way to verify attainment of the stopping payoff. Reflection requires the corresponding Neumann or oblique derivative condition and a well-posed reflected state process. Unbounded-state applications require growth bounds and localization. These conditions are checked separately from a neural residual.

For a smooth function $v$, write
\begin{align}
 \LL^a v&=b(t,x,a)^\top Dv+\tfrac12\tr\{\sigma(t,x,a)\sigma(t,x,a)^\top D^2v\},\label{eq:generator}\\
 \HH^a[v]&=f(t,x,a,v)+\LL^a v,\qquad
 F[v]=-v_t-\sup_{a\in\Act(t,x)}\HH^a[v].\label{eq:hjb}
\end{align}
The value solves $F[V^*]=0$ with the specified boundary data, in the viscosity sense where classical derivatives do not exist. The diffusion covariance is the positive semidefinite matrix $\sigma\sigma^\top$. No claim that the product of two arbitrary positive semidefinite matrices is symmetric positive semidefinite is needed.

A continuous martingale representation can also be used when its quadratic variation is Markov in an augmented state. If $d\langle M\rangle_s=Q(s,S_s,a_s)ds$, the covariance in \eqref{eq:generator} becomes $\sigma Q\sigma^\top$. A merely predictable process $Q_s$ is not, without further state augmentation, a deterministic coefficient of a Markov HJB equation. Quadratic variation by itself also does not imply conditionally Gaussian increments. The implementation in this paper uses the Brownian model \eqref{eq:state}.

\subsection{Infinite horizon and economic normalization}
For a stationary problem, the boundary-value equation is $\sup_a\{f(x,a,V)+\LL^aV\}=0$. We use an infinite-horizon result only after specifying admissibility, integrability, and the selected utility process. In the time-additive case, a sufficient verification requirement is that the localized It\^o identities are uniformly integrable and $\E[e^{-\rho T}|V(S_T)|]\to0$ for the policies being compared. Recursive problems require an appropriate utility comparison and transversality condition instead of simply borrowing the additive one.

Cardinal transformations also require care. Multiplying a fixed utility function by a positive constant does not change its maximizer. Adding a function of an \emph{endogenously controlled preference state} generally does. Consumption units, preference normalization, and terminal utility are therefore primitive economic choices in Section~\ref{sec:ndu}.

% R15 component; only input paths and enumerated location/grammar prose are edited.
% Source: revisions/2026-10-04-r15/manuscript/method.tex; commit 1cb3cc9efe135966ec228dfaf51c6841a6dfc97c.
% Generated by code/integrate.py; edit extraction rules there.
% Reviewed source: ECTA.tex, line 111, commit f5021cefa71492babfcbfa580e0c984f59a9de26
\section{Neural Bellman Operators}\label{sec:method}
The actor $a_\phi(t,x)$ and critic $v_\xi(t,x)$ represent policies and continuation values. An internal actor and an external actor can be separate parameter blocks; their joint action remains the input to the same economic generator. A smooth activation is useful for exact differentiation, while a nonsmooth neural representation can be used with the positive-weight Bellman formulation below.

\subsection{Feasibility and boundary data}
For a box $[\underline a,\bar a]$, a sigmoid map produces interior feasible actions but does not attain its endpoints at finite weights. Binding constraints can instead be represented by a projection onto the closed interval, by an explicit boundary candidate, or by a finite action set including the endpoints. The constrained portfolio calculation uses a projection and can attain a zero risky share exactly. On a simplex, a probability map supplies a feasible interior point and boundary faces must again be considered when certifying the optimum. General constraints require a feasible parameterization, a certified projection, or a constrained action solver. A finite penalty is not an exact feasibility certificate.

On a finite horizon, an elementary hard-terminal architecture is
\begin{equation}\label{eq:terminal}
 v_\xi(t,x)=G(x)+(T-t)N_\xi(t,x).
\end{equation}
This enforces the terminal face, but not automatically a lateral boundary. A compatible lifting $G_B$ and a factor $d_B$ vanishing on all Dirichlet faces give $v_\xi=G_B+d_BN_\xi$. Alternatively, boundary residuals are trained separately and retained explicitly in the error bound. A reflected preference state requires its own derivative condition; a terminal architecture does not enforce it. For recursive utility, a domain-preserving transformation must additionally keep the continuation value in the domain of the aggregator.

\subsection{The block-specific computational graph}
Let $\nu$ be a specified distribution of state--time points and let $\sg$ denote stop-gradient. A critic step minimizes
\begin{equation}\label{eq:criticloss}
 L_C(\xi;\bar\phi)=\E_\nu\left[-\partial_t v_\xi-
                 \HH^{\sg(a_{\bar\phi})}[v_\xi]\right]^2+L_B(\xi),
\end{equation}
where $L_B$ is absent only for boundary conditions enforced exactly. The policy, sampling locations, and any rival actions are detached. The value, time derivative, state gradient, and Hessian remain differentiable with respect to the critic parameters.

An actor step minimizes
\begin{equation}\label{eq:actorloss}
 L_A(\phi;\bar\xi)=-\E_\nu\HH^{a_\phi}
 [\sg(v_{\bar\xi}),\sg(Dv_{\bar\xi}),\sg(D^2v_{\bar\xi})].
\end{equation}
All critic derivative channels are fixed in this step. There is no critic gradient from $L_A$ and no actor gradient from $L_C$. A printed sum $L_A+\omega L_C$ may be a reporting device, but it is not jointly differentiated. The implementation exposes separate functions and regression tests for these derivative paths.

The distinction is substantive: joint differentiation can move
the critic away from the policy-evaluation root. Appendix~\ref{app:r15evaluation}
gives a counterexample, the recursive lifetime policy derivative, and the
conditioning requirements. The local Hamiltonian objective is not identified
with a lifetime policy gradient under arbitrary state sampling.
The evaluated reference policy may remain fixed while a separate candidate
actor is improved at successive work levels. Updating that reference defines
a new evaluation target and is recorded explicitly.

\subsection{Rollout and costate evaluation}\label{sec:r11learning}
% Reviewed source: revisions/2026-10-04-r11/manuscript/learning.tex, line 4, commit f5021cefa71492babfcbfa580e0c984f59a9de26
For a frozen policy $a^{(j)}$, a positive-time Euler rollout from $(t,y)$ supplies a sample return $\widehat V^{(j)}(t,y)$ and its pathwise derivative $\widehat p^{(j)}(t,y)$. Differentiation includes the state dependence of subsequent policy actions. Freezing the policy means freezing its parameters, not deleting that derivative. The targets evaluate the current policy, not the optimal policy, and contain no maximizing action labels. For fixed finite weights, bounded actions and smooth bounded production justify pathwise differentiation of this finite rollout away from the harmless clipping kinks. These are derivatives of the discretized evaluation problem, not asserted unbiased derivatives of its continuous-time limit.

The critic minimizes
\begin{equation}\label{eq:r11critic}
 \widehat{\E}\left[
  0.05\{v_\theta-\widehat V^{(j)}\}^2+
  \|d(D_yv_\theta-\widehat p^{(j)})\|_d^2\right].
\end{equation}
The factor $d$ puts normalized marginal values on an order-one scale. The rollout implementation uses two hidden tanh layers and a hard terminal component. Holding its new derivative fixed, the actor then maximizes
\begin{equation}\label{eq:r11actor}
 \widehat{\E}\left[
 \overline{\log a_\phi}-\tfrac\zeta2\bar a_\phi^2
                          -a_\phi\mathbin{\cdot}\operatorname{sg}(D_yv_\theta)\right].
\end{equation}
Terms independent of actions have been omitted. Consumption is represented directly by $m_0(t)+\epsilon\tanh N_\phi(t,y)$; there is no online Hamiltonian search in its deployment. Value evaluation, costate evaluation, and feasible improvement are distinct operations, with their own stored losses and work counts. The number of critic and actor steps per rollout batch is a declared design parameter. Failure of either numerical block is recorded, not interpreted as successful policy improvement.

The stochastic gradient targets make the economically relevant marginal continuation values visible to the evaluation loss. This is not a new Bellman principle, nor an assumption that derivative regression invariably dominates differential residual fitting. It is a concrete implementation of the NBO evaluation map. The preserved differential implementation and its two independent trace banks continue to test a different evaluation procedure. Both feed the same economic distinction between policy evaluation and feasible improvement.

% Reviewed source: revisions/2026-09-29-r6/manuscript/operator.tex, line 1, commit f5021cefa71492babfcbfa580e0c984f59a9de26
\subsection{The implemented map and its verification output}\label{sec:implemented}
Fix the economy, boundary realization, network classes, optimization budget, and verification procedure, denoted collectively by $\vartheta$. Let $\mathcal P$ be the admissible feedback policies, $\Xi$ the critic parameter space, and $\Omega$ the sampling streams. The implemented NBO is the numerical map
\begin{equation}\label{eq:implemented-r6}
 \mathcal N_\vartheta:\mathcal P\times\Xi\times\Omega
 \longrightarrow\mathcal V_\vartheta\times\mathcal P_\vartheta\times\mathcal C_\vartheta,
 \qquad(a,\xi,\omega)\longmapsto(v,\widetilde a,C).
\end{equation}
The first component is approximate policy evaluation; the second is feasible improvement; the third records verification of the resulting pair. A record identifies the model, parameters, domain, boundary faces, evaluation error, maximization error, coverage, discretization errors, and unchecked assumptions. An unavailable bound is not zero. Failure to obtain a useful certificate is a possible output of the map. With its sampling stream fixed, the implementation is deterministic. We do not assert that a network learns an unrestricted Bellman operator on a function space.

The evaluation--improvement structure is an instance of approximate policy iteration, including its partial-evaluation variants \citep{scherrer2015}. Neural representation does not introduce a new principle of optimality, and monotonicity belongs to the underlying positive-weight scheme, not to gradient descent. The NBO specification joins the block-specific derivative graph to boundary and utility-domain restrictions and to a common economic error account. Sections~\ref{sec:nonquadratic} and~\ref{sec:ndu-neural} implement the differential and positive-weight versions of this map, respectively, with independent verification of their outputs.

% R15 component; only input paths and enumerated location/grammar prose are edited.
% Source: revisions/2026-10-04-r15/manuscript/algorithm.tex; commit 1cb3cc9efe135966ec228dfaf51c6841a6dfc97c.
\begin{algorithm}[t]
\caption{Neural Bellman evaluation, feasible improvement, and economic assessment}\label{alg:nbo}
\begin{algorithmic}[1]
\Require Economic primitives and information; feasible actions and utility domain; boundary data; evaluation variant; training, validation and stopping rules; independent final assessment.
\State Declare the feasible reference policy to be evaluated. Initialize a feasible candidate and a boundary-compatible continuation representation.
\For{each permitted evaluation--improvement iteration}
 \State Freeze the declared reference policy, rivals and evaluation target. Fit only the critic by the declared differential, positive-weight, or rollout rule. The reference may remain fixed across work levels.
 \State Freeze the critic and its value/derivative channels. Update only the feasible actor, or apply the declared feasible action solver.
 \State Record all evaluation and improvement work, corrections, unsuccessful checks, and the resulting deployment rule.
 \State Apply the predetermined validation and stopping rule, using no final-test outcomes.
\EndFor
\State Freeze the selected policy and continuation object before independent final evaluation.
\State Compute the specified economic endpoints and direct contrasts, with their approximation, implementation and statistical accounts.
\State Return the fitted objects, their implementation and domain, and the complete accuracy and work record.
\end{algorithmic}
\end{algorithm}

\section{Economic Decisions and Accuracy Targets}\label{sec:r19targets}
A task $\theta$ and observed state $s$ determine a feasible action set
$A_\theta(s)$ and current reward $r_\theta(s,a)$. The postdecision state is
$T_\theta(s,a)$. Write
\begin{equation}
 Q_\theta(s,a)=r_\theta(s,a)+S_\theta\{T_\theta(s,a)\},\qquad
 L_\theta(s,a)=\sup_{b\in A_\theta(s)}Q_\theta(s,b)-Q_\theta(s,a).
 \label{eq:r19objective}
\end{equation}
The continuation $S_\theta$ includes future rewards under a specified future
policy, conditional on the postdecision information. Current rewards are
excluded. A change in the current charge alone need not change $S_\theta$;
a change in the future policy generally does. For an observed catalogue
$\{(s_i,\theta_i)\}_{i=1}^q$, our new primary target is
$q^{-1}\sum_i L_{\theta_i}(s_i,a_i)\leq\varepsilon$.

\begin{table}[htbp]\centering\small
\caption{Economic targets and their verification domains}\label{tab:r19targets}
\begin{tabular}{p{.23\textwidth}p{.34\textwidth}p{.34\textwidth}}
\toprule Target & Comparison object & Interpretation in this paper\\\midrule
Continuation risk & Exact conditional continuation at specified arguments &
Prediction diagnostic; neither policy welfare nor control regret.\\
Finite-menu loss & Best action in one exogenous finite action menu &
Decision diagnostic at common tasks; not a continuous-action bound.\\
Catalogue regret & Best of sixteen historical method--budget candidates &
Earlier simultaneous payoff intervals and their exact closure.\\
Scalar-interval regret & True maximum over $[.05,\bar a_i]$ in the stored finite economy &
New primary accuracy certificate and realized stopping target.\\
Diffusion control regret & Full declared adapted-control class in continuous time &
Separate boundary, implementation, diffusion and information accounts.\\
Utility or game criterion & Policy utility, spike deviation or unilateral best response &
Application-specific admissibility and verification conditions.\\\bottomrule
\end{tabular}
\end{table}

The word ``regret'' below always refers to the displayed comparison class.
The current scalar experiment takes $s\in\R^d$ but restricts the current
withdrawal to a uniform scalar action. Neither a large state dimension nor
full enumeration of its finite innovation law turns this target into an
optimum over vector actions or all adapted future controls.

\section{From a Learned Continuation to Economic Accuracy}\label{sec:r19decision}
\subsection{Centered approximation error and the selected decision}
Fix an anchor continuation $S_0$ and its approximation $\widehat S_0$.
For each task, let $e=\widehat S_0-S_0$ and
$\Delta_\theta=S_\theta-S_0$. The implemented action $\widehat a$ satisfies
\begin{equation}
 r_\theta(s,\widehat a)+\widehat S_0(T_\theta(s,\widehat a))
 \geq \sup_{a\in A_\theta(s)}
 \{r_\theta(s,a)+\widehat S_0(T_\theta(s,a))\}-\delta_\theta(s).
 \label{eq:r19greedy}
\end{equation}
Here $\delta_\theta(s)\geq0$ is an actual optimization allowance, not a
training-loss tolerance. For a function on the feasible actions, write
$\operatorname{osc}_A f=\sup_A f-\inf_A f$.

\begin{proposition}[Continuation error, transport, and decision loss]
\label{prop:r19decision}
Suppose the suprema in \eqref{eq:r19objective} and
\eqref{eq:r19greedy} are finite. Then
\begin{equation}
 L_\theta(s,\widehat a)\leq\delta_\theta(s)
  +\operatorname{osc}_{A_\theta(s)}(e\circ T_\theta)
  +\operatorname{osc}_{A_\theta(s)}(\Delta_\theta\circ T_\theta).
 \label{eq:r19osc}
\end{equation}
For a finite action set $\{a_1,\ldots,a_m\}$ with $m\geq2$ and fixed
positive weights $\pi_j$ summing to one, define
\begin{align*}
 R_\pi(s)&=\sum_j\pi_j(e(T_\theta(s,a_j))-\bar e_\pi(s))^2,\\
 \bar e_\pi(s)&=\sum_j\pi_j e(T_\theta(s,a_j)),
 &C_\pi&=\max_{i\ne j}(\pi_i^{-1}+\pi_j^{-1}).
\end{align*}
If $|\Delta_\theta|\leq b_\theta$ on the queried postdecision set, then
\begin{equation}
 L_\theta(s,\widehat a)\leq\delta_\theta(s)
                       +\sqrt{C_\pi R_\pi(s)}+2b_\theta.
 \label{eq:r19riskloss}
\end{equation}
Both inequalities hold for each fitted object and each realized cache.
Expectations over training or cache randomness may subsequently be taken.
\end{proposition}

The proof is in Section~\ref{app:r19proofs}. An additive continuation error
common to all feasible actions has no effect on the decision. This is why
an absolute squared prediction error may be large even when decisions are
accurate. Conversely, the risk at only the NBO/reference action pair cannot
control the losses of another method choosing elsewhere. Our new diagnostic
uses a common exogenous action set and, separately, the arguments selected
by each predictor. Inequality~\eqref{eq:r19riskloss} is a bound for the common
finite set; it is not applied to the entire scalar interval using just three
sampled actions.

If the true finite-menu best action has gap $g>0$ from every competitor,
then a bound strictly smaller than $g$ implies that the implemented action
is best. Without such a margin, arbitrarily small centered prediction error
can change the identity of the selected action. Neither statement orders two
methods' realized payoffs from their risks alone.

\subsection{A nontrivial reuse region when the future changes}
Let $X_{k+1}^\theta=F_{\theta,k}(X_k^\theta,Z_k)$ describe the future
under the policy being evaluated. A common-innovation coupling need not
make the transition maps identical. On an invariant state domain, assume
\begin{align}
 \|F_{0,k}(x,z)-F_{0,k}(x',z)\|&\leq L_k\|x-x'\|,
 &\|F_{\theta,k}(x,z)-F_{0,k}(x,z)\|&\leq\epsilon_{\theta,k},\notag\\
 |g_{0,k}(x)-g_{0,k}(x')|&\leq H_k\|x-x'\|,
 &|g_{\theta,k}(x)-g_{0,k}(x)|&\leq\eta_{\theta,k}.
 \label{eq:r19transportassumptions}
\end{align}
These bounds concern future rewards and transitions, not just the current
task. Include the terminal payoff as $g_{\theta,K}$ and use nonnegative
reward weights $\omega_k$. The initial postdecision state is shared.

\begin{theorem}[Continuation transport and refresh]
\label{thm:r19transport}
Set $d_0=0$ and
$d_{k+1}=L_kd_k+\epsilon_{\theta,k}$. Under
\eqref{eq:r19transportassumptions}, uniformly on the common postdecision
domain,
\begin{equation}
 |S_\theta-S_0|\leq
 b_\theta:=\sum_{k=0}^{K}\omega_k
                   (\eta_{\theta,k}+H_kd_k).
 \label{eq:r19transport}
\end{equation}
Thus the anchor may be reused at a task whenever the right side of
\eqref{eq:r19osc}, or the applicable finite-menu bound
\eqref{eq:r19riskloss}, is at most the desired economic tolerance.

More generally, suppose a recursive evaluation on a common invariant utility
domain has maps $V_{\theta,k}=\mathcal G_{\theta,k}(V_{\theta,k+1})$.
If these maps are sup-norm Lipschitz with constants $\beta_k$ on that domain,
$\|\mathcal G_{\theta,k}(v)-\mathcal G_{0,k}(v)\|_\infty\leq\nu_{\theta,k}$,
and $\|V_{\theta,K}-V_{0,K}\|_\infty\leq b_K$, then
\begin{equation}
 \|V_{\theta,0}-V_{0,0}\|_\infty\leq
 \sum_{k=0}^{K-1}\left(\prod_{j=0}^{k-1}\beta_j\right)\nu_{\theta,k}
       +\left(\prod_{j=0}^{K-1}\beta_j\right)b_K.
 \label{eq:r19recursive}
\end{equation}
Empty products equal one.
\end{theorem}

The theorem supplies a class of tasks with genuinely different futures.
For example, uniformly bounded transition and reward perturbations of an
anchor economy satisfy it on any verified invariant compact domain with
bounded Lipschitz constants. A single learned continuation then supports a
whole neighborhood of counterfactuals until its economic error allowance is
spent. The neighborhood is conditional on the displayed stability bounds;
it is not a distribution-free assertion about all future policies. When
future rewards, constraints, or opponents' policies violate those bounds,
the evaluated continuation must be recomputed or its transport allowance
reestablished. For recursive utility the common utility domain is essential:
a local Lipschitz constant cannot be extended across a singular utility
boundary. For games, this comparison evaluates a specified deviation problem;
it does not replace the fixed-point or best-response conditions in the
Economic Applications companion.

\subsection{An explicit amortization class}
The preceding inequalities give a sufficient class in which a shared scalar
continuation supports a cheaper declared accuracy procedure, not just a
post hoc interpretation of successful cells. Fix a finite exogenous query
family with $K$ distinct postdecision arguments and a common continuation
$S(x)=\phi(x)^{\mathsf T}\beta$. The feature map $\phi:\mathcal X\to\R^p$
may be a frozen neural dictionary or a conventional basis. Its construction
cost is charged, and it is fixed before the fitting noises below. At design
points $x_i$, observe $Y_i=S(x_i)+\xi_i$, where the noises are conditionally
independent, mean zero, and $\sigma$-sub-Gaussian. Let $\Phi$ have rows
$\phi(x_i)^{\mathsf T}$ and assume
\[
 G_n=\Phi^{\mathsf T}\Phi\text{ is invertible},\qquad
 \max_{x\text{ queried}}\phi(x)^{\mathsf T}G_n^{-1}\phi(x)\leq\ell/n.
\]
For example, $G_n/n\succeq cI$ and $\|\phi(x)\|^2\leq B_\phi^2$
imply $\ell=B_\phi^2/c$. These are checkable design conditions, not a
claim about an arbitrary trained representation.

\begin{corollary}[A sufficient learned-continuation work advantage]
\label{cor:r19amortization}
Fix $\varepsilon>\delta\geq0$ and $0<\alpha<1$. Least-squares continuation
fitting and $\delta$-accurate finite-menu maximization have loss at most
$\varepsilon$ on every query with probability at least $1-\alpha/2$ if
\begin{equation}
 n\geq \frac{8\sigma^2\ell}{(\varepsilon-\delta)^2}
                    \log\frac{4K}{\alpha},
 \label{eq:r19learnbudget}
\end{equation}
provided the displayed full-rank and leverage conditions hold. Direct
sample-average predictions of the same continuation have the same guarantee
when they use at least
$t=\lceil8\sigma^2\log(4K/\alpha)/(\varepsilon-\delta)^2\rceil$
observations per distinct argument, under the analogous scalar noise bound.
Common innovations across arguments are allowed.

Suppose the complete declared learned procedure costs
$F(n)+Kc_N+C(q)$ and the declared direct procedure costs
$C_0(t)+Kt c_R+C(q)$, with all fitting, caching, selection and verification
included. Both procedures meet the target jointly with probability at least
$1-\alpha$, and the learned procedure is strictly cheaper whenever
\begin{equation}
 K(t c_R-c_N)>F(n)-C_0(t).
 \label{eq:r19explicitadvantage}
\end{equation}
Repeated identical arguments are counted once in $K$. If verification or
selection costs differ, their actual difference enters the right side.
\end{corollary}

The class contains nonlinear continuations represented by finite smooth
neural features, not only linear economic models. The leverage bound keeps
fitting accuracy uniform as queries are reused. A bounded representation
bias or transported future adds its explicit allowance to the decision bound
and reduces the available margin in \eqref{eq:r19learnbudget}. The result
compares two fully specified sufficient-accuracy procedures; it is not a
lower bound on the best possible simulation method. A Raw procedure may
itself fit a surrogate, and then belongs to the learned comparison class.
In particular, this corollary neither certifies the assumptions for the
experiment's jointly trained network nor asserts that its conventional
surrogates must be slower. It provides an analytical reuse class that the
numerical results need not be asked to establish by extrapolation.

\subsection{Actual stopping and the amortization threshold}
Consider a feasible scalar action in $[\ell,u]$ and a twice differentiable
true objective with $Q''\leq-\kappa<0$. Define the directional residual
\[
 r(a)=\begin{cases}
 \max\{Q'(a),0\},&a=\ell,\\
 \max\{-Q'(a),0\},&a=u,\\
 |Q'(a)|,&\ell<a<u.
 \end{cases}
\]
An interval enclosure of $Q'(a)$ and a lower bound on $\kappa$ give a
computable upper bound $U(a)$ in the following statement.
\begin{proposition}[Certified stopping and complete work]\label{prop:r19stopping}
For every feasible scalar action,
\begin{equation}
 \sup_{b\in[\ell,u]}Q(b)-Q(a)\leq\frac{r(a)^2}{2\kappa}.
 \label{eq:r19residual}
\end{equation}
For a finite sequence of candidate-fitting stages, let a service stop at the
first stage whose mean outward bound is at most $\varepsilon$. If each bound
is a deterministic enclosure for the true task objectives, the returned
service satisfies the target regardless of the preceding failed checks.
If checks instead have conditional error allocations $\alpha_j$ given the
past and $\sum_j\alpha_j\leq\alpha$, the analogous conclusion holds with
probability at least $1-\alpha$.

The realized cost of a service stopping at stage $J$ is the sum of all
initialization, fitting or cache construction, querying, verification, and
durable-record costs through $J$. In an exact fixed-plus-linear cost model
$W_m(q)=F_m+q c_m$, where $c_m$ includes marginal verification and all
attempted stages, two methods that both meet the same target satisfy
\begin{equation}
 W_N(q)<W_R(q)\quad\Longleftrightarrow\quad
 q(c_R-c_N)>F_N-F_R.
 \label{eq:r19break}
\end{equation}
When $c_R>c_N$, this gives the strict break-even threshold. If only cost
intervals are available, a guaranteed ranking requires an upper bound for
one method below a lower bound for the other, not two upper bounds.
\end{proposition}

The cost model is an interpretation, not an imposed shape for the measured
frontier. Fitting stages may change with query volume, and verification can
dominate marginal prediction work. We therefore execute every volume rather
than extrapolate a line from a single 384-query comparison. The new experiment
runs each method's frozen stage sequence; it does not search globally for a
winner and then claim to have paid only that winner's bill. All comparison
services and shared startup are reported separately from the service that
an economist would deploy after choosing a method.

\section{An All-State Policy Certificate}\label{sec:r23policy}
The sup-norm account above is natural on bounded value domains. In a capital
model with unbounded shocks, a uniform bound on an absolute local advantage
can instead be infinite. A cost-weighted account resolves this difficulty
without restricting verification to states visited by the fitted policy.
We state it for additive evaluation. The recursive-utility and strategic
models retain the domain and deviation conditions in
Section~\ref{sec:r21composition} and the Economic Applications companion.

Use a cost-minimization convention in this section: a reward is the negative
of the corresponding cost. Let $\ell_t(x,a)\geq0$, $h(x)\geq0$, and
$0<\beta\leq1$. For an admissible policy $\pi$, define
\begin{equation}
 U_t(x)=J_t^\pi(x)=\E_x^\pi\left[
 \sum_{s=t}^{T-1}\beta^{s-t}\ell_s(X_s,A_s)+\beta^{T-t}h(X_T)\right].
 \label{eq:r23ownvalue}
\end{equation}
The comparison class includes every adapted policy with finite cost. Assume
that the conditional expectations below exist for these policies. The
candidate's local improvement opportunity is
\begin{equation}
 g_t(x)=U_t(x)-\inf_a\{\ell_t(x,a)+\beta\E[U_{t+1}(X_{t+1})\mid x,a]\}.
 \label{eq:r23advantage}
\end{equation}
In particular, $g_t$ uses the candidate's actual future, not an optimal
continuation or a fitted proxy for it.

\begin{theorem}[Coercive policy certificate]\label{thm:r23coercive}
Suppose that $U_T=h$ and, for one finite $\eta\geq0$,
\begin{equation}
 0\leq g_t(x)\leq\eta\ell_t(x,a)
 \quad\text{for every date, state, and admissible action.}
 \label{eq:r23coercivity}
\end{equation}
Then, with $J_t^*$ the infimum over all adapted finite-cost policies,
\begin{equation}
 0\leq U_0(x)-J_0^*(x)
 \leq\frac{\eta}{1+\eta}U_0(x).
 \label{eq:r23relative}
\end{equation}
The statement is deterministic when its premises are deterministic. A
sampled certificate must establish all its premises on a common event.
\end{theorem}

\begin{proof}
For any admissible action, the definition of $g_t$ and
\eqref{eq:r23coercivity} imply
$U_t(x)\leq(1+\eta)\ell_t(x,a)+\beta\E[U_{t+1}\mid x,a]$.
Apply this inequality along an arbitrary adapted comparison policy and
iterate conditional expectations. At the terminal date,
\[
 U_0(x)\leq(1+\eta)\E\sum_{t<T}\beta^t\ell_t
                 +\beta^T\E h\leq(1+\eta)J_0(x),
\]
where the last step uses $h\geq0$. Taking the infimum over comparison
policies and rearranging proves the assertion. The lower bound follows
because the candidate is admissible.
\end{proof}

The argument is a Bellman-residual comparison, not a new principle of
optimality. Its useful feature is that a bound proportional to a coercive
stage cost can hold on an unbounded domain even when an absolute sup-norm
bound cannot. It is not enough to check \eqref{eq:r23coercivity} on a finite
catalogue. Nonnegative costs and the comparison-policy integrability
requirements are substantive premises.

\subsection{A multisector capital-adjustment economy}
Let $x\in\R^d$ denote capital stocks relative to a reference path and
$u\in\R^d$ their investment adjustments. There are $T$ reoptimized dates:
\begin{align}
 X_{t+1}&=A_tX_t+B_tu_t+\xi_{t+1},&
 \xi_{t+1}&\sim N(0,\Sigma),\label{eq:r23dynamics}\\
 \ell_t(x,u)&=x'Q_tx+u'R_tu,&h(x)&=x'Q_Tx.\label{eq:r23cost}
\end{align}
Shocks are independent across dates and correlated across sectors.
$Q_t$ and $R_t$ are positive definite; $Q_T$ is positive semidefinite.
The problem is posed in discrete time, on the full state and action spaces.
It is a quadratic tracking economy, not a calibrated quantitative description
of a particular industry and not a discretization claim about the original
diffusion model. Its continuous innovation support and vector controls make
it a direct test of the policy certificate, rather than a larger scalar
catalogue.

For a stored linear policy $u_t=-K_tx$, its \emph{own} value is
$U_t(x)=x'P_tx+c_t$, where, writing $F_t=A_t-B_tK_t$,
\begin{align}
 P_T&=Q_T,&c_T&=0,\nonumber\\
 P_t&=Q_t+K_t'R_tK_t+\beta F_t'P_{t+1}F_t,&
 c_t&=\beta\{c_{t+1}+\tr(P_{t+1}\Sigma)\}.
 \label{eq:r23evaluation}
\end{align}
This is evaluation of a supplied policy. It contains no optimization over
future gains. Define
\begin{equation}
 H_t=R_t+\beta B_t'P_{t+1}B_t,\qquad
 D_t=H_tK_t-\beta B_t'P_{t+1}A_t.
 \label{eq:r23residual}
\end{equation}
Completing the square in the full vector action gives the exact identity
$g_t(x)=x'D_t'H_t^{-1}D_tx$. In particular, an action mesh is unnecessary.

\begin{corollary}[Uniform capital-policy bound]\label{cor:r23capital}
Let $q_t\leq\lambda_{\min}(Q_t)$ and
$r_t\leq\lambda_{\min}(R_t)$ be strictly positive lower bounds. Set
\begin{equation}
 \eta=\max_{t<T}\frac{\|D_t\|_F^2}{q_tr_t}.
 \label{eq:r23eta}
\end{equation}
For every state, the stored policy satisfies \eqref{eq:r23relative}.
Consequently, on $\|x\|_2^2\leq b$, its policy loss is at most
\begin{equation}
 \varepsilon_{\rm pol}
 =\frac{\eta}{1+\eta}\{b\|P_0\|_\infty+c_0\}.
 \label{eq:r23ball}
\end{equation}
The comparison is with all adapted finite-cost policies and all real vector
actions, not only linear policies or the class used to construct the critic.
\end{corollary}

\begin{proof}
$H_t\succeq R_t\succeq r_t I$, so
$g_t(x)\leq\|D_t\|_F^2\|x\|_2^2/r_t
\leq\eta x'Q_tx\leq\eta\ell_t(x,u)$ for every $u$.
Theorem~\ref{thm:r23coercive} applies. The true $P_0$ is symmetric positive
semidefinite; hence $x'P_0x\leq\|x\|_2^2\|P_0\|_\infty$.
Finite-cost comparison policies have finite second moments at every relevant
date by coercivity, which supplies the required integrability.
\end{proof}

The coefficient matrices, rather than sampled state outcomes, establish
coverage. Thus the class-approximation allowance is zero: the verifier already
compares with the unrestricted adapted class. The value-transfer allowance
is also zero for this declared discrete-time economy. Neither zero is
transferred to a diffusion approximation, a constrained nonlinear economy,
or an equilibrium calculation.

\subsection{Arithmetic and implementation}
The verifier encloses \eqref{eq:r23evaluation}--\eqref{eq:r23ball} using
outward matrix-ball arithmetic. Stored binary coefficients define the model
and candidate exactly. Each matrix product encloses rounding and propagated
input uncertainty; the code rejects invalid coercivity constants, overflow,
and nonfinite entries. It does not use a computed optimal policy. The
supplement states the product enclosure and the independent tests.

A separate implementation account allows a state-dependent action error
$e_t$ satisfying $\|e_t(x)\|_2\leq\epsilon_{\rm exec}\|x\|_2$ at every
date. A backward quadratic upper envelope supplies
$\varepsilon_{\rm impl}$, so the final stopping target is
\begin{equation}
 \varepsilon_{\rm pol}+\varepsilon_{\rm impl}\leq\varepsilon.
 \label{eq:r23stop}
\end{equation}
This is a guarantee for the declared relative-error class, not an assumption
that a physical floating-point executor is overflow-free on every real input.
The exact stored policy corresponds to zero execution error. The experiment
also reports the outward bound for $\epsilon_{\rm exec}=10^{-12}$.

\subsection{From a trained neural factor to policy accuracy}
The preceding verifier accepts any candidate. The next result concerns the
\emph{trained representation} itself. For the square-activation continuation
$\widehat U_t(x)=\|W_tx\|_2^2/d+c_t$, every entry of $W_t$ is trainable.
This is not a fixed-feature regression assumption. Define its population
own-policy derivative loss by
\[
 \mathcal L_t(W_t;P_t)=\tfrac14\|W_t'W_t-dP_t\|_F^2,
 \qquad G_t=\nabla_{W_t}\mathcal L_t=W_t(W_t'W_t-dP_t).
\]
Here $P_t$ must be evaluated under the \emph{returned} policy, even if the
last training target came from an earlier policy. Let
$\widehat P_t=W_t'W_t/d$ and $\widehat P_T=Q_T$. The residual of the actual
actor relative to its fitted continuation is
\[
 Z_t=(R_t+\beta B_t'\widehat P_{t+1}B_t)K_t
                   -\beta B_t'\widehat P_{t+1}A_t.
\]

\begin{proposition}[Trainable-factor certificate]\label{prop:r23factor}
Suppose $\sigma_{\min}(W_{t+1})\geq s_{t+1}>0$. Put
$a_{t+1}=\|G_{t+1}\|_F/(ds_{t+1})$ and $a_T=0$. Then
\begin{equation}
 \|D_t\|_F\leq\|Z_t\|_F+
       \beta\|B_t\|_2\|F_t\|_2 a_{t+1}.
 \label{eq:r23factor}
\end{equation}
Replacing $\|D_t\|_F$ in \eqref{eq:r23eta} by these upper bounds gives a
valid full-policy certificate. A failed rank bound or a large gradient
leaves this particular certificate unresolved.
\end{proposition}
\begin{proof}
Invertibility gives
$\|\widehat P_{t+1}-P_{t+1}\|_F
\leq\|G_{t+1}\|_F/(ds_{t+1})$. Direct subtraction shows
$D_t=Z_t+\beta B_t'(\widehat P_{t+1}-P_{t+1})F_t$.
The matrix-norm inequality proves \eqref{eq:r23factor}; the preceding
corollary then applies.
\end{proof}

This proposition supplies an a posteriori implication for a jointly trained
hidden factor, rather than asserting that a nonconvex optimizer reaches
small loss or that every neural architecture has this property. The factor
and a conventional quadratic span the same quadratic value class in this
experiment. Whether training the factor is worth its additional work is
therefore an empirical comparison, not a conclusion of the proposition.
The reported factor diagnostic re-evaluates the returned policy, encloses
its gradient and rank bound, and retains any unresolved diagnostic. It is a
deterministic analysis of the frozen candidates, not a replacement stopping
rule selected after observing their performance.

\section{A full-policy certificate for recursive risk preferences}
\label{sec:r27entropic}
The order-preserving composition argument does not by itself supply a useful
uniform Lipschitz constant for an exponential aggregator on unbounded states.
We give a separate, executable bridge for that case. Keep the capital dynamics
and vector actions in \eqref{eq:r23dynamics}, but evaluate a policy recursively by
\begin{equation}\label{eq:r27entropicbellman}
 U_t(x)=\ell_t(x,-K_tx)+\beta\mathcal R_\theta U_{t+1}(F_tx),
 \qquad \mathcal R_\theta v(y)=\theta^{-1}\log\E e^{\theta v(y+\xi)},
\end{equation}
where $F_t=A_t-B_tK_t$, $\theta>0$, and $\xi\sim N(0,\Sigma)$.
The limit at $\theta=0$ is conditional expectation. This is a recursive
entropic criterion with discounting outside the certainty equivalent, not an
assertion that it equals every exponential-of-total-cost convention.
Exponential-quadratic control and its structural solutions are classical
\citep{r27jacobson1973,r27cui2023}. We use that structure to verify a returned
neural policy without using an optimal value as a training label or certificate.

For a quadratic continuation $x'Px+c$, Gaussian integration gives
\begin{align}
 \mathcal R_\theta v(y)&=y'\Psi_\theta(P)y+c+h_\theta(P),\notag\\
 \Psi_\theta(P)&=P(I-2\theta\Sigma P)^{-1},\qquad
 h_\theta(P)=-\frac{\log\det(I-2\theta\Sigma P)}{2\theta}.
 \label{eq:r27gaussian}
\end{align}
The sufficient domain condition $2\theta\|\Sigma\|_2\|P\|_2<1$
is checked at every date; a failed domain check is not silently truncated.
Own-policy evaluation yields $U_t(x)=x'P_tx+c_t$. Define
\[
 H_t=R_t+\beta B_t'\Psi_\theta(P_{t+1})B_t,\qquad
 D_t=H_tK_t-\beta B_t'\Psi_\theta(P_{t+1})A_t.
\]
Choose bounds $\zeta_t\geq\|D_t\|_F$, $f_t\geq\|F_t\|_2$,
$b_t\geq\|B_t\|_2$, $p_{t+1}\geq\|P_{t+1}\|_2$,
$\sigma\geq\|\Sigma\|_2$, and $\tau\geq\tr\Sigma$.
With $\delta_T=\kappa_T=0$, compute backward
\begin{align}
 \gamma_t&=1-2\theta\sigma p_{t+1},&
 w_t&=\delta_{t+1}/\gamma_t^2,\notag\\
 \bar f_t&=f_t+\frac{b_t}{r_t}
             (\zeta_t+\beta b_tf_tw_t),\notag\\
 \delta_t&=\frac{\zeta_t^2}{r_t}+\beta\bar f_t^2w_t,&
 \kappa_t&=\beta\left(\kappa_{t+1}+
                   \frac{\tau\delta_{t+1}}{\gamma_t}\right).
 \label{eq:r27riskenvelope}
\end{align}

\begin{theorem}[Recursive full-policy comparison]\label{thm:r27entropic}
Suppose $Q_t\succeq q_tI\succ0$, $R_t\succeq r_tI\succ0$, $Q_f\succeq0$,
and $\Sigma\succeq0$. If every $\gamma_t$ is positive and
$\delta_t\leq q_t$ for $1\leq t<T$, then
\begin{equation}\label{eq:r27riskgap}
 0\leq U_0(x)-J_0^*(x)\leq\delta_0\|x\|^2+\kappa_0
 \qquad(x\in\R^d).
\end{equation}
Here $J^*$ compares all adapted policies with finite recursive cost, not only
linear policies, and the controls range over all of $\R^d$.
\end{theorem}

The certificate is a Bellman subsolution: subtract
$\delta_t\|x\|^2+\kappa_t$ from the policy's own value. The matrix and constant
allowances are different because the Gaussian log determinant contributes a
state-independent cost. The derivative of $\Psi_\theta$ supplies
$\gamma_t^{-2}$, while that of $h_\theta$ supplies $\gamma_t^{-1}$.
These checked quantities replace an uninstantiated global utility Lipschitz
constant. The proof and the separately computed execution allowance are in
Section~\ref{sec:r27riskproof}.

The verifier encloses own-policy matrices with outward matrix balls. Inverses
are proposed numerically and accepted only after a Neumann-residual test.
It does not require evaluation of an optimal policy, an enumeration of the
Gaussian law, or a state catalogue. In the declared discrete model, class
approximation and value-transfer allowances are zero because the comparison
is direct; this does not set those allowances to zero for a discretized
continuous economy or for the separate strategic models.

For the neural construction, a square-activation critic learns
$d\Psi_\theta(P_{t+1})$ after the future policy has been finalized. A separately
run structural recursion uses the same analytic information without the
factor fit. The comparison therefore isolates the additional neural training
cost rather than concealing the structural solution. Neither Gaussian
integration nor the classical structural optimum is claimed as a new method.
The new conclusion concerns the verified propagation of a returned critic's
and policy's errors through a nonlinear recursive preference map.

\section{Training to a Recursive Policy Target}\label{sec:r31constructive}
The policy certificate and the training procedure answer different questions.
A certificate assesses a returned policy. A constructive training statement
must determine an admissible budget before observing that policy's evaluation
target. This section joins the two statements for the trainable square critic
in the recursive capital economy of Section~\ref{sec:r27entropic}. The original
smooth-network and controlled-diffusion formulations remain unchanged.

\subsection{Evaluation without interval expansion}
Let $C_T=Q_f$ and let $C_t$ be arbitrary symmetric numerical proposals for the
coefficients of the returned policy's own value. They need not be computed by
an optimal-control recursion. Write $F_t=A_t-B_tK_t$ and suppose that the
following residual and matrix norms are enclosed from above:
\[
 \rho_t\geq\|C_t-Q_t-K_t'R_tK_t-
                  \beta F_t'\Psi_\theta(C_{t+1})F_t\|_2.
\]
Starting at $e_T=0$, define
\begin{align}
 p_{t+1}&=\|C_{t+1}\|_2+e_{t+1},&
 \gamma_t&=1-2\theta\sigma p_{t+1},\notag\\
 e_t&=\rho_t+\beta\|F_t\|_2^2e_{t+1}/\gamma_t^2,
 \label{eq:r31tube}
\end{align}
where $\sigma\geq\|\Sigma\|_2$. Norms in a numerical implementation mean
certified upper bounds, and the computed $\gamma_t$ is a lower bound.

\begin{proposition}[Residual-centered evaluation]\label{prop:r31tube}
If every margin in \eqref{eq:r31tube} is positive, the returned policy's
recursive value is finite and its exact coefficients satisfy
$\|P_t-C_t\|_2\leq e_t$. Moreover, a valid bound on its full-action residual is
\begin{equation}
 \zeta_t=
 \|R_tK_t-\beta B_t'\Psi_\theta(C_{t+1})F_t\|_F
 +\beta\|B_t\|_2\|F_t\|_F e_{t+1}/\gamma_t^2.
 \label{eq:r31defect}
\end{equation}
These quantities can be inserted directly into
Theorem~\ref{thm:r27entropic} and its implementation account.
\end{proposition}

The result is independent of the proposal algorithm. An inaccurate proposal
increases its checked Bellman residual; it does not become a valid evaluation
merely because it is labeled as one. Propagating a scalar spectral radius
around each newly evaluated center avoids feeding an expanding componentwise
box through successive matrix inverses. This changes neither the economic
criterion nor the stopping tolerance. The certificate remains uniform over
real states and full vector deviations, rather than over a stored catalogue.

\subsection{A budget chosen from economic primitives}
Fix a reference feedback $L_t$ and a positive slack $\eta$. This feedback is a
feasibility witness, not an optimal solution or a training label. Let
$Q_t\succeq q_tI\succ0$, $R_t\succeq r_tI\succ0$, and set
$q_* =\min_{t<T}q_t$. Construct a scalar coefficient envelope backward:
\begin{align}
 \bar p_T&\geq\|Q_f\|_2,&
 \bar\gamma_t&=1-2\theta\sigma\bar p_{t+1}>0,\notag\\
 \bar p_t&\geq\|Q_t+L_t'R_tL_t+\eta Q_t\|_2
 +\beta\|A_t-B_tL_t\|_2^2\bar p_{t+1}/\bar\gamma_t.
 \label{eq:r31primitive}
\end{align}
The entire envelope is computed before any policy target or critic fit.
For $t<T-1$, put
\[
 m_t=dq_{t+1},\qquad M_t^{\rm ub}=d\bar p_{t+1}/\bar\gamma_t,
 \qquad
 f_t=\|A_t\|_2\left(1+
        \frac{\beta\|B_t\|_2^2M_t^{\rm ub}}{dr_t}\right).
\]
Use the same expression for $f_{T-1}$ with the terminal coefficient bound.
Let $b_t\geq\|B_t\|_2$ and $\tau\geq\operatorname{tr}\Sigma$. Define
\begin{align}
 \widetilde f_t&=f_t+\frac{b_t}{r_t}
       \left(\sqrt{\eta r_tq_t}
                +\frac{\beta b_tf_tq_*}{\bar\gamma_t^2}\right),\notag\\
 \lambda_T&=\mu_T=0,\qquad
 \lambda_t=q_t\mathbf1_{\{t<T-1\}}
           +\frac{\beta\widetilde f_t^2}{\bar\gamma_t^2}\lambda_{t+1},
 \notag\\
 \mu_t&=\beta\left(\mu_{t+1}
                +\frac{\tau\lambda_{t+1}}{\bar\gamma_t}\right).
 \label{eq:r31weights}
\end{align}
For initial states satisfying $\|x_0\|^2\leq b_0$, write
$W_0=b_0\lambda_0+\mu_0$. Choose $s>0$ such that
\begin{equation}
 s\leq\eta,\qquad s\lambda_t\leq q_*\ (1\leq t<T),\qquad
 sW_0\leq\varepsilon_{\rm train}.
 \label{eq:r31allocate}
\end{equation}
Zero weights impose no restriction. The remaining allowance
$\varepsilon-\varepsilon_{\rm train}$ is reserved for numerical implementation;
it is not presumed sufficient until independently checked.

At date $t<T-1$, the finalized future policy supplies the own-policy target
$M_t=d\Psi_\theta(P_{t+1})$. Train the scalar continuation
$\widehat U_{t+1}(x)=\|W_tx\|^2/d+\widehat c_{t+1}$ by
\begin{equation}
 W^{(j+1)}=W^{(j)}-\alpha_tW^{(j)}
          \big((W^{(j)})'W^{(j)}-M_t\big),\quad
 W^{(0)}=\sqrt{c_t}\,O_t,
 \label{eq:r31training}
\end{equation}
where $O_t$ is orthogonal, $0<c_t\leq m_t$ and
$0<\alpha_t\leq1/(4M_t^{\rm ub})$. All entries of $W$ are updated.
For $\beta b_tf_t>0$, the required Gram error is
\begin{equation}
 a_t=\frac{d\sqrt{s r_tq_t}}{\beta b_tf_t},\qquad
 \|(W^{(j)})'W^{(j)}-M_t\|_F\leq a_t.
 \label{eq:r31trainingtarget}
\end{equation}
A continuation-independent action needs no such fit. The terminal actor uses the known transformed coefficient
$\Psi_\theta(Q_f)$ without a factor fit, as in \eqref{eq:r32terminal}. Every actor minimizes its fitted full quadratic action
objective, and its own future is then evaluated for the preceding date.

\begin{theorem}[Primitive-budget neural policy construction]
\label{thm:r31constructive}
Suppose the primitive envelope \eqref{eq:r31primitive} exists with positive
margins. In exact arithmetic, the backward construction just specified
terminates after finitely many hidden-weight updates and returns a policy
with
\begin{equation}
 0\leq J_0^\pi(x)-J_0^*(x)
       \leq s\{\lambda_0\|x\|^2+\mu_0\}
       \leq\varepsilon_{\rm train},\qquad \|x\|^2\leq b_0.
 \label{eq:r31constructivegap}
\end{equation}
It suffices to take, at each fitted date,
\begin{equation}
 k_t=\left\lceil
 \frac{\log^+(\sqrt d\,M_t^{\rm ub}/a_t)}
      {-\log(1-\alpha_tc_t)}\right\rceil
 \label{eq:r31cap}
\end{equation}
updates. Neither an optimal value nor an already-fitted target is an input
to the envelope, tolerance, step size, initialization scale, or cap.
The comparator in \eqref{eq:r31constructivegap} is the full adapted
finite-recursive-cost class in the declared economy.
\end{theorem}

Section~\ref{sec:r31proofs} proves the theorem by joint backward induction on
the policy envelope and the critic target. This removes a circular use of a
successful fitted policy to justify its own training budget. The conclusion
is architecture-specific: it concerns a trainable square factor with a
controlled initialization, not a frozen feature dictionary, arbitrary
nonconvex initialization, or every network used in the original applications.
The class admits an inexpensive structural solution; that comparator remains
part of the numerical experiment. Finite neural trainability is not itself
a claim that neural training is the least expensive way to solve this class.

The implemented procedure uses conservative caps, retains every failed
threshold, and checks the rounded returned policy by
Proposition~\ref{prop:r31tube}. Its final acceptance rule is the original
full-policy tolerance, including execution error. Thus the constructive
real-arithmetic theorem does not silently promote a rounded training trace
to an economic certificate.

\section{Warm-Start Learning to a Full-Policy Accuracy Target}
\label{sec:r34robust}
A local training guarantee is useful for policy construction only when its
error allowance is paid in the policy objective. There is a further difficulty
with warm starts: their fitted Gram matrices need not underestimate the new
continuation. The primitive envelope used for the cold construction cannot
therefore be carried over by an appeal to monotonicity. We give an explicit
budget that accommodates overshooting factors, errors in own-policy
evaluation, and inexact actor solves. The economic model and the comparison
with all adapted policies are those of Section~\ref{sec:r27entropic}.

\subsection{An implementable local acceptance condition}
At date $t$, write $S_t=\Psi_\theta(P_{t+1}^{\pi})$, where the future policy
has already been finalized. Suppose an evaluator supplies a symmetric center
$\widetilde M_t$ and a certified radius
\[
 \|\widetilde M_t-dS_t\|_2\leq e_t^M.
\]
The fitted scalar continuation has coefficient
$\widehat S_t=W_t'W_t/d$; $W_t$ may be rectangular. Put
\begin{align}
 F_t&=A_t-B_tK_t,&
 Z_t&=R_tK_t-\beta B_t'\widehat S_tF_t,\notag\\
 a_t&\geq\|W_t'W_t-\widetilde M_t\|_F,&
 z_t&\geq\|Z_t\|_F.\label{eq:r34localdata}
\end{align}
Thus $z_t$ measures the actor's equation error, not its fitting objective.
The exact full-action residual is
$D_t=R_tK_t-\beta B_t'S_tF_t$. For any certified bounds
$b_t\geq\|B_t\|_2$ and $f_t\geq\|F_t\|_2$,
\begin{equation}\label{eq:r34mixed}
 \|D_t\|_F\leq z_t+
 \frac{\beta b_tf_t}{d}(a_t+\sqrt d\,e_t^M).
\end{equation}
The two matrix errors enter differently. The Gram error is bounded in
Frobenius norm, whereas the evaluation radius is spectral. Omitting the
factor $\sqrt d$ from the latter term would not be valid. When an enclosure
of $\|F_t\|_F$ is available, it replaces $\sqrt d f_t$ in that term.

\subsection{A primitive budget that allows overshooting}
Keep the reference feedback, slack $\eta>0$, coefficient envelope
$\bar p_t$, and positive margins $\bar\gamma_t$ from
\eqref{eq:r31primitive}. Let
\[
 q_* = \min_{t<T}q_t,\qquad
 u_t=\bar p_{t+1}/\bar\gamma_t,\qquad
 v_t=\sqrt{\eta r_tq_t}.
\]
All quantities in this paragraph are computed before fitting any target.
Define the enlarged closed-loop bound
\begin{equation}\label{eq:r34f}
 f_t^+=\|A_t\|_2
       \left\{1+\frac{\beta b_t^2(u_t+q_*)}{r_t}\right\}
             +\frac{b_tv_t}{r_t},
 \qquad
 g_t=f_t^++\frac{b_t}{r_t}
       \left(v_t+\frac{\beta b_tf_t^+q_*}{\bar\gamma_t^2}\right).
\end{equation}
Norms in a numerical budget denote certified upper bounds. Starting from
$\lambda_T^+=\mu_T^+=0$, compute
\begin{align}
 \lambda_t^+&=q_t+
       \frac{\beta g_t^2}{\bar\gamma_t^2}\lambda_{t+1}^+,
 &\mu_t^+&=\beta\left(\mu_{t+1}^+
             +\frac{\tau\lambda_{t+1}^+}{\bar\gamma_t}\right).
 \label{eq:r34weights}
\end{align}
Unlike the exact-terminal construction, this account permits an actor error
at the last date. For an initial radius $\|x_0\|^2\leq b_0$, choose $s>0$
so that
\begin{equation}\label{eq:r34allocation}
 s\leq\eta,\qquad s\lambda_t^+\leq q_*\ (1\leq t<T),\qquad
 s(b_0\lambda_0^++\mu_0^+)\leq\varepsilon_{\rm policy}.
\end{equation}
The local acceptance rule is
\begin{equation}\label{eq:r34accept}
 a_t+e_t^M\leq dq_*,\qquad
 z_t+\frac{\beta b_tf_t^+}{d}(a_t+\sqrt d\,e_t^M)
                         \leq\sqrt{s r_tq_t}.
\end{equation}
Both inequalities are checked. The first controls the candidate's magnitude;
the second pays for its decision error. A positive training tolerance remains
only after the evaluator and actor have taken their shares of the budget.

\begin{theorem}[Budgeted warm-start policy construction]
\label{thm:r34robust}
Suppose the primitive coefficient envelope has positive moment margins and
\eqref{eq:r34allocation} holds. Construct the policy backward. At every date,
evaluate the finalized future policy, return a positive-semidefinite fitted
coefficient and an actor, and verify \eqref{eq:r34localdata}--\eqref{eq:r34accept}.
Then $0\preceq P_t^\pi\preceq\bar p_tI$ and
\begin{equation}\label{eq:r34policybound}
 0\leq J_0^\pi(x)-J_0^*(x)
 \leq s\{\lambda_0^+\|x\|^2+\mu_0^+\}
 \leq\varepsilon_{\rm policy},\qquad \|x\|^2\leq b_0.
\end{equation}
The comparison is with all adapted policies of finite recursive cost, over
all real vector controls. No commutation or one-sided approximation condition
on the fitted Gram matrices is required.

For a warm fit whose center, initial factor, step error, and precision satisfy
Theorem~\ref{thm:r33warm}, a finite cap follows by substituting its positive
remaining Gram allowance in \eqref{eq:r33cap}. The evaluator's radius and the
actor equation error must still satisfy \eqref{eq:r34accept}. A training gate
alone is not the acceptance rule for the economic policy.
\end{theorem}

The proof in Section~\ref{sec:r34proof} jointly establishes the coefficient
envelope and the full-policy bound. The extra $q_*$ in \eqref{eq:r34f} pays
for approximation on either side of the target; the final term pays for the
actor solve. The construction does not use an optimal policy as a training
label. At the last date the true continuation coefficient is
$\Psi_\theta(Q_f)$, not $Q_f$ unless $\theta=0$ (or the transform is otherwise
unchanged). The terminal target may be supplied directly and its actor
checked without training another continuation.

\begin{corollary}[A finite arithmetic implementation]\label{cor:r34finite}
Suppose the economic matrices, risk parameter, discount factor, and supplied
bounds are rational, and the stated moment margins are strict. Own-policy
coefficient evaluation and actor equations can be solved by exact rational
matrix arithmetic. At a certified warm start, compute the factor update
exactly and round its entries to a dyadic grid. Choosing the grid fine enough
that its Gram noise floor is below the remaining positive allowance gives a
finite number of updates satisfying \eqref{eq:r34accept}. When a warm gate
fails, the controlled cold construction with rational squared initialization
provides a finite alternative at nonterminal dates. Consequently the policy
bound has a finite-arithmetic realization, rather than only a formal
real-arithmetic interpretation.
\end{corollary}

This corollary is a statement about a specified arithmetic routine. Ordinary
floating-point matrix products require their own full-step error enclosure;
nearest-storage rounding by itself does not cover them. The independent
rounded-policy verifier is retained for the larger numerical experiments.
An additional action-execution allowance is added to
\eqref{eq:r34policybound} before comparing with a requested tolerance.

\subsection{Construction, reliability, and work}
For the declared admissible inputs, the deterministic cap is uniform over
warm factors and permitted update errors. It is not a reliability interval
estimated from a few fitting seeds. Its economic content is a finite
construction to a full-policy tolerance under recursive preferences,
including the cost of precision. The exact-arithmetic example in the
replication package checks the policy envelope, noncommuting Gram errors,
evaluation, and the returned vector actors against the same rational inputs.
It is a constructive validation instance, not an additional independent
sample in the earlier economic experiments.

An amortization comparison includes evaluating the changed future, testing
the warm gate, unsuccessful checks, fallback training, actor solves,
verification, and durable output. The theory gives a cap for the neural
factor updates; it does not turn that cap into a lower wall-clock cost than
an analytic structural solution. The earlier cost frontiers and their
unfavorable comparisons are retained. In particular, improving the trainable
continuation's mathematical guarantee does not change which procedure was
cheapest in an already executed experiment.

\section{Precision-Adaptive Neural Continuation Refresh}
\label{sec:r37precision}
An old continuation can be inaccurate for a changed economy and nevertheless
provide a useful initialization for learning its new continuation. A useful
refresh rule must specify both how much learning is required and how accurately
that learning must be implemented. We give such a rule for the trainable
square continuation, allowing rectangular factors and noncommuting targets.
The economic stopping criterion remains the full-policy bound of
Theorem~\ref{thm:r34robust}.

Fix one date after the future policy has been finalized. During this fit the
symmetric target $M$ is fixed and satisfies $mI\preceq M\preceq LI$, with
$0<m\leq L$. For $W_j\in\R^{p\times d}$, $p\geq d$, let
\[
 C_j=M^{-1/2}W_j'W_jM^{-1/2},\qquad E_j=C_j-I.
\]
The implemented update is
\begin{equation}\label{eq:r37update}
 W_{j+1}=\tfrac12 W_j\{3I-M^{-1}W_j'W_j\}+\Delta_j,
 \qquad \|\Delta_jM^{-1/2}\|_2\leq\nu_j.
\end{equation}
Every entry of the hidden factor is trainable. The inverse is that of the
current fixed target; it is neither an inverse of the optimal value nor a
matrix transported unchanged from a different economic future. In whitened
coordinates, the ideal update is a Newton--Schulz polar iteration. Matrix
polar iterations and their numerical analysis are classical
\citep{r37higham2008}. We use their error identity to allocate precision within
an own-policy economic error budget, rather than claim a new matrix iteration.

Choose a checked basin radius $0<r\leq1/2$ and put $c_r=(3+r)/4$. If
$\|E_0\|_2\leq\rho_0\leq r$, define $\rho_{j+1}=\rho_j^2$. For every
step that is needed, require
\begin{equation}\label{eq:r37precision}
 2\nu_j+\nu_j^2\leq(1-c_r)\rho_j^2.
\end{equation}
A zero initial error needs no update. The condition budgets the whole
implemented step. It cannot be justified by bounding only the final storage
conversion when preceding matrix products have also been rounded.

\begin{theorem}[Precision-adaptive refresh]\label{thm:r37precision}
Suppose the fixed target, initial basin, and full-step errors satisfy the
conditions just stated. Then every returned factor obeys
\begin{equation}\label{eq:r37double}
 (1-\rho_j)M\preceq W_j'W_j\preceq(1+\rho_j)M,
 \qquad \rho_j=\rho_0^{2^j}.
\end{equation}
For a requested relative tolerance $0<\tau<1$ with $\tau<\rho_0$, a sufficient number of updates is
\begin{equation}\label{eq:r37cap}
 K=\left\lceil\log_2\left\{
       \frac{\log(1/\tau)}{\log(1/\rho_0)}\right\}\right\rceil.
\end{equation}
The number is zero if $\rho_0\leq\tau$. No simultaneous eigenbasis for
$W_0'W_0$ and $M$ is assumed.

Suppose, more specifically, that matrix products and target solves are exact
and each of the $pd$ entries of the updated factor is rounded to the nearest
multiple of $2^{-b_j}$. Let $A$ be any finite certified upper bound for
$\sqrt{pd/m}$. Choosing the first positive integer $b_j$ for which
\begin{equation}\label{eq:r37bits}
 2\nu_j+\nu_j^2\leq(1-c_r)\rho_j^2,
 \qquad \nu_j=A2^{-b_j-1},
\end{equation}
is sufficient. Such an integer exists at every required step. For a constant
$C$ depending only on $A$ and $r$, the chosen integers satisfy
\begin{equation}\label{eq:r37bitaccount}
 \sum_{j=0}^{K-1}b_j\leq CK+4\log_2(1/\tau).
\end{equation}
A checked inverse of the fixed target may be constructed once and applied
at each update. Thus a fit requiring updates needs one target-inverse
construction and $K$ target-inverse applications; a zero-update fit needs
neither.
\end{theorem}

The proof is in Section~\ref{sec:r37proof}. The first conclusion is a
uniform deterministic statement over the declared basin and perturbations.
The second is a fractional-storage account, not a bound on all bit operations
in exact rational arithmetic. It omits neither inverse applications nor
residual and positivity checks from a measured service cost. Numerators and
denominators in exact intermediate products can be considerably larger than
the stored factor precision. An implementation with a fixed precision cap
must fail explicitly if it cannot satisfy \eqref{eq:r37bits}; the theorem
does not authorize it to return an uncertified factor.

\subsection{A target change and its economic allowance}
Suppose an old factor has relative error at most $\rho_{\rm old}$ for
$M_{\rm old}$, and evaluation of the new future establishes
\[
 (1-\omega)M_{\rm new}\preceq M_{\rm old}
                \preceq(1+\omega)M_{\rm new},\qquad 0\leq\omega<1.
\]
Then the same factor has relative error at most
$\rho_{\rm old}+\omega+\rho_{\rm old}\omega$ for $M_{\rm new}$. If
this upper bound does not exceed $r$, it supplies an admissible initial
radius for Theorem~\ref{thm:r37precision}. Otherwise, a direct basin test or
a controlled cold construction is required. This is a condition for using
an old factor as an initialization, not for accepting its old policy without
improvement. The old inverse is not reused across the target change.

\begin{corollary}[Precision-budgeted full-policy construction]
\label{cor:r37policy}
Under the primitive envelope and allocation of
Theorem~\ref{thm:r34robust}, let each positive remaining Gram allowance be
$a_t$, and let the independently evaluated center satisfy
$m_tI\preceq\widetilde M_t\preceq L_tI$. A basin-certified refresh with
\begin{equation}\label{eq:r37economic}
 \tau_t\leq\min\{1/2,a_t/(\sqrt d\,L_t)\}
\end{equation}
reaches that Gram allowance after the finite schedule in
\eqref{eq:r37cap}--\eqref{eq:r37bits}. If the evaluation radius and returned
actor also satisfy \eqref{eq:r34accept}, then
\[
 0\leq J_0^\pi(x)-J_0^*(x)
       \leq s\{\lambda_0^+\|x\|^2+\mu_0^+\}
       \leq\varepsilon_{\rm policy},\qquad \|x\|^2\leq b_0.
\]
The comparison is with all adapted policies of finite recursive cost and all
real vector controls in the stated economy. A separate action-execution
allowance is added when the implemented actions differ from the stored
rational feedback.
\end{corollary}

Thus the economic primitives determine the envelope and policy allocation;
own-policy evaluation supplies the current target and its radius; the
precision rule constructs a trainable continuation to the remaining
allowance; and the actor and full-policy checks complete the certificate.
The terminal target is the known transformed coefficient
$\Psi_\theta(Q_f)$, not $Q_f$ in general. A direct terminal evaluation
requires no neural factor fit.

The update count and fractional precision account concern a specified neural
construction. They do not rank all ways of solving the economy. In particular,
the structural solution legitimately bypasses factor training. The numerical
comparison below charges both constructions to the same economic stopping
target, includes that competitor, and reports the work of evaluating, checking,
and writing the result as well as the work of training.

\section{Constructive Neural Bellman Operators beyond quadratic continuations}
\label{sec:r39nonlinear}

A continuation method has an economic use when its approximation can be
translated into a decision guarantee without assuming that the approximation
class contains the value function. We now give this translation for nonlinear
recursive economies. The comparison is with all feasible adapted policies.
The construction uses compact state and action sets, not a Gaussian coefficient
identity. The preceding diffusion framework and unbounded-state results remain
in force under their respective assumptions; compactness is imposed only in
this additional realization.

\subsection{A residual-oscillation theorem}
Let $X$ and $A$ be compact metric spaces, let $F_t:X\times A\times Z\to X$
and $c_t:X\times A\to\R$ be continuous, and let $Z$ have a finite distribution
with fixed probabilities $q_z$. Write $\beta_t>0$ for the date-$t$ continuation
weight, $B_0=1$, and $B_t=\prod_{j<t}\beta_j$. The conditional continuation
functional $\rho_t$ is monotone and cash invariant:
$\rho_t(Y+a)=\rho_t(Y)+a$. It is continuous on finite vectors. In particular,
these conditions hold for
\begin{equation}\label{eq:r39risk}
 \rho_\theta(Y)=\theta^{-1}\log\sum_zq_z e^{\theta Y(z)},\qquad
 \rho_0(Y)=\sum_zq_zY(z),\qquad \theta\geq0.
\end{equation}
Boundedness makes the exponential moment finite. Define
\[
 Q_t(v;x,a)=c_t(x,a)+\beta_t\rho_t(v(F_t(x,a,Z))),\qquad
 T_tv(x)=\min_{a\in A}Q_t(v;x,a).
\]
Continuity gives measurable minimizers. Recursive costs for arbitrary adapted
policies use the same conditional functional at every realized history. Let
$V_T^*=g$ and $V_t^*=T_tV_{t+1}^*$.

\begin{theorem}[Nonlinear full-policy comparison]\label{thm:r39comparison}
Let $v_t$ be continuous functions and let $\pi_t(x)\in A$ be measurable.
Suppose that, at every state,
\begin{align}
 \ell_t\leq v_t-T_tv_{t+1}\leq u_t,&\qquad t<T,\label{eq:r39residual}\\
 Q_t(v_{t+1};x,\pi_t(x))-T_tv_{t+1}(x)&\leq\eta_t,\label{eq:r39actor}\\
 \ell_T\leq v_T-g&\leq u_T.\nonumber
\end{align}
Then, for every initial state and against all feasible adapted policies,
\begin{equation}\label{eq:r39loss}
 0\leq J_0^\pi-V_0^*
 \leq\sum_{t<T} B_t(u_t-\ell_t+\eta_t)+B_T(u_T-\ell_T).
\end{equation}
Moreover, $v_0-\sum_{t\leq T}B_tu_t$ is a lower bound on $V_0^*$.
Neither membership of $V^*$ in a neural class nor convergence of a training
optimizer is required for this comparison.
\end{theorem}

\begin{proof}
Monotonicity and cash invariance imply
$\|T_tv-T_tw\|_\infty\leq\beta_t\|v-w\|_\infty$; the same statement holds
for the operator of any fixed feasible action. Backward induction applied to
\eqref{eq:r39residual} yields
\[
 v_t-\sum_{s=t}^T(B_s/B_t)u_s\leq V_t^*
 \leq v_t-\sum_{s=t}^T(B_s/B_t)\ell_s.
\]
The terminal equality is immediate. For the induction step, shift the next
function by the corresponding deterministic constant, apply $T_t$, and use
$T_tv_{t+1}\in[v_t-u_t,v_t-\ell_t]$.
For the returned policy, \eqref{eq:r39actor} instead gives
\[
 J_t^\pi\leq v_t+\sum_{s=t}^{T-1}(B_s/B_t)(\eta_s-\ell_s)
                    -(B_T/B_t)\ell_T.
\]
Subtract the lower bound on $V_t^*$. At an arbitrary history, each realized
action has continuation cost at least the Bellman minimum, and induction
compares its future cost with $V_{t+1}^*$. Thus the Bellman value is a lower
bound for every adapted policy, not only for the Markov or neural class.
A measurable minimizing Markov policy attains it. This proves both claims.
\end{proof}

The economically relevant residual is its oscillation $u_t-\ell_t$.
A common additive valuation error cancels from the loss comparison. This is
the nonlinear counterpart of centered continuation error in the earlier
operator analysis, rather than an assumption that a critic learns absolute
value levels accurately.

\subsection{A finite construction and a trained-network check}
Take $X=[0,1]$ and $A=[0,\bar a]$. Suppose $c_t$ and $F_t$ have known
Lipschitz constants separately in state and action. If $v_{t+1}$ has
Lipschitz constant $L_{t+1}$, admissible constants for $Q_t$ are
\begin{equation}\label{eq:r39lipschitz}
 L_{x,t}=L_{c,x,t}+\beta_t L_{F,x,t}L_{t+1},\qquad
 L_{a,t}=L_{c,a,t}+\beta_t L_{F,a,t}L_{t+1}.
\end{equation}
The same $L_{x,t}$ bounds the state variation of $T_tv_{t+1}$.
Use a uniform state partition $x_i=ih$ and an action grid of mesh $\Delta$,
including both endpoints. At every pair, verified arithmetic supplies
$\underline q_{ij}\leq Q_t(v_{t+1};x_i,a_j)\leq\overline q_{ij}$.
Set $L_i=\min_j\underline q_{ij}$, $U_i=\min_j\overline q_{ij}$ and
$\delta_t=L_{a,t}\Delta/2$. Let $v_t$ be the linear interpolant of stored
nodal values $y_i$, with no requirement that these equal an exact minimizer.
Then valid residual constants are
\begin{align}\label{eq:r39mesh}
 \ell_t&=\min_i(y_i-U_i)-L_{x,t}h/2,\\
 u_t&=\max_i(y_i-L_i+\delta_t)+L_{x,t}h/2.\nonumber
\end{align}
If the action $a_{j_i}$ is stored at node $i$ and extended by the
nearest-node rule, a valid actor allowance is
\begin{equation}\label{eq:r39meshactor}
 \eta_t=\max_i(\overline q_{i j_i}-L_i)+\delta_t+L_{x,t}h.
\end{equation}
All expressions are enclosed outward in the implementation.

\begin{proposition}[Constructive nonquadratic continuation]\label{prop:r39construction}
Under the preceding conditions, \eqref{eq:r39mesh}--\eqref{eq:r39meshactor}
and a certified terminal interpolation error imply
Theorem~\ref{thm:r39comparison}. Every such $v_t$ has the exact representation
\begin{equation}\label{eq:r39relu}
 v_t(x)=y_0+s_0x+\sum_{i=1}^{N-1}(s_i-s_{i-1})(x-x_i)_+,
 \qquad s_i=(y_{i+1}-y_i)/h.
\end{equation}
For each fixed finite horizon, a feasible neural policy with any prescribed
positive loss tolerance can be constructed in finitely many operations if
arithmetic precision can be increased to deliver the required interval widths.
This assertion concerns the constructive mesh algorithm, not unrestricted
stochastic-gradient training.
\end{proposition}

\begin{proof}
The continuous action minimum at a node lies in $[L_i-\delta_t,U_i]$:
a grid point is within $\Delta/2$ of any minimizing action. On a state cell,
linear interpolation of nodal Bellman values differs from the Bellman function
by at most $L_{x,t}h/2$, since the interpolation-weighted distance to the two
endpoints is $2r(h-r)/h\leq h/2$. This proves \eqref{eq:r39mesh}.
Moving from a state to its nearest node changes both its chosen-action value
and the Bellman minimum by at most $L_{x,t}h/2$ each, proving
\eqref{eq:r39meshactor}. Slope changes prove \eqref{eq:r39relu}.

For finiteness, choose nodal midpoints of $[L_i,U_i]$ and the action minimizing
$\overline q_{ij}$. Suppose each $Q$ interval has width at most $\omega$ and
each stored midpoint has an additional rounding error at most $\zeta$.
Then $U_i-L_i\leq\omega$. Residual oscillation plus actor allowance is at most
\[
 2\omega+2\zeta+2\delta_t+2L_{x,t}h.
\]
For the Lipschitz recursion, each nodal midpoint is within
$\delta_t+\omega+\zeta$ of the exact Bellman value. Consequently the next
interpolant has Lipschitz constant at most
$L_{x,t}+2(\delta_t+\omega+\zeta)/h$.
Choose $\Delta\leq h^2$ and $\omega+\zeta\leq h^2$ at every date.
Equations~\eqref{eq:r39lipschitz} then give, by finite backward induction,
a bound on all interpolant Lipschitz constants uniform over $0<h\leq1$.
Terminal interpolation contributes at most $L_g h/2+\zeta$ in absolute error.
The right-hand side of \eqref{eq:r39loss} therefore tends to zero with $h$.
A sufficiently fine finite grid and sufficiently precise finite arithmetic
attain any positive target. This proof does not assert that fixed binary64
arithmetic suffices for every target.
\end{proof}

A conventional spline routine and \eqref{eq:r39relu} represent the same
constructed function; relabeling it alone is not a neural work advantage.
The second implementation trains all affine hidden weights in
$f_t(x)=a_t+b_tx+\sum_{j=1}^p c_{tj}(w_{tj}x+d_{tj})_+$.
It uses the certified nodal values as deterministic training data and selects
actions from its own continuation predictions. Its selected actions are then
checked by \eqref{eq:r39meshactor} against the independent interval brackets.
The full cost of the reference construction, training and verification is
charged. Acceptance certifies the returned policy even when the optimizer has
no convergence guarantee. A rejected proposal can be replaced by the
constructive policy above; this gives a terminating certified procedure without
asserting that every trained proposal must be accepted.

\begin{proposition}[All-state verification of trained ReLU continuations]
\label{prop:r39breakpoints}
Let $f$ be the preceding one-hidden-layer scalar ReLU network and $v$ a linear
spline. The maximum of $|f-v|$ on $[0,1]$ occurs at an endpoint, a spline knot,
or a neural root $-d_j/w_j$ in that interval, with $w_j\ne0$. Enclosing each
root by outward division and evaluating both functions over its interval gives
a certified uniform bound, including nondyadic roots. Direct evaluation costs
$O((N+p)(p+\log N))$ arithmetic operations; sorting the breakpoints, when used,
costs $O((N+p)\log(N+p))$.
\end{proposition}
\begin{proof}
Between consecutive points in the specified set, $f-v$ is affine. The absolute
value of an affine function is convex and is bounded by its endpoint maximum.
Continuity includes the endpoints of adjacent pieces. Zero hidden slopes
produce no breakpoint. Every exact breakpoint belongs to its outward root
interval, so interval evaluation bounds its value even when a root cannot be
stored exactly. A direct network evaluation uses $O(p)$ operations and a
binary-search spline evaluation $O(\log N)$. There are at most $N+p+1$ points.
\end{proof}

This certificate removes the mesh-Lipschitz remainder used in the earlier
network-to-spline check. It does not retroactively change the earlier policy
bounds or execution clocks. The new endpoint audit uses exact rational
arithmetic only as an independent check of the verifier, not as the production
method used in the floating-point economic services.

\subsection{Nonlinear investment and the value of reoptimization}
\label{sec:r39economy}
The economic realization has bounded capital, irreversible investment and a
shortfall penalty. Its primitives are
\begin{align}
 F(x,a,z)&=\operatorname{clip}_{[0,1]}
       \{13x/16+a+x(1-x)/16+z\},\label{eq:r39capital}\\
 c_P(x,a)&=(x-11/16)^2+Pa^2+4a^4+2(3/8-x)_+^2,\nonumber\\
 g(x)&=2(x-11/16)^2+2(3/8-x)_+^2,\qquad \beta=15/16,\nonumber
\end{align}
where $a\in[0,1/4]$ and the shocks $(-1/16,0,1/16)$ have probabilities
$(1/4,1/2,1/4)$. The price $P$ measures the quadratic investment burden;
$\theta$ in \eqref{eq:r39risk} measures aversion to high future costs.
A low-capital household faces both a distance-to-target loss and a shortfall
penalty. Investment is capped and cannot be negative. The quartic adjustment
cost, nonlinear accumulation and clipping prevent a Riccati reduction.
This is a theoretical counterfactual economy, not an estimated calibration.

Valid primitive constants are $L_{c,x}=23/8$, $L_{F,x}=7/8$,
$L_{c,a}=P/2+1/4$, $L_{F,a}=1$, and $L_g=17/4$.
They follow by differentiation on smooth pieces and continuity at the kinks;
clipping is nonexpansive. Hence Proposition~\ref{prop:r39construction}
verifies its assumptions directly from the primitives.

\begin{proposition}[Fixed-rule investment exposure]\label{prop:r39exposure}
For a fixed feasible policy rule $\pi$ in \eqref{eq:r39capital}, independent of
$P$, the derivative $D_t^\pi=\partial J_t^\pi/\partial P$ exists along each
finite shock tree and satisfies $D_T^\pi=0$ and
\begin{equation}\label{eq:r39exposure}
 D_t^\pi=a_t^2+\beta\sum_z\widetilde q_{t,z}D_{t+1}^\pi(F_t),\qquad
 \widetilde q_{t,z}=
 \frac{q_z\exp(\theta J_{t+1}^\pi(F_t))}
      {\sum_yq_y\exp(\theta J_{t+1}^\pi(F_t))}.
\end{equation}
At $\theta=0$, $\widetilde q=q$. In particular $D_t^\pi\geq0$.
For old and new rules, the total cost change decomposes exactly as
\begin{align}\label{eq:r39decomposition}
 J_0^{\pi^{\rm new}}(4)-J_0^{\pi^{\rm old}}(1)
 ={}&\{J_0^{\pi^{\rm old}}(4)-J_0^{\pi^{\rm old}}(1)\}\\
 &-\{J_0^{\pi^{\rm old}}(4)-J_0^{\pi^{\rm new}}(4)\}.\nonumber
\end{align}
The first bracket is mechanical repricing; the second is the cost saved by
reoptimization at the new price.
\end{proposition}
\begin{proof}
For a fixed rule, the states and selected actions on the finite shock tree do
not depend on $P$, even if the rule is discontinuous in state. Each stage cost
is affine in $P$. Differentiate the finite log-sum-exp recursion; its derivative
is the probability-weighted derivative in \eqref{eq:r39exposure}. Backward
induction proves existence and nonnegativity. Add and subtract
$J_0^{\pi^{\rm old}}(4)$ to obtain \eqref{eq:r39decomposition}.
\end{proof}

No global monotonicity of the optimal investment rule is inferred from this
fixed-rule statement. The numerical counterfactual below instead reports the
actual policy changes and an independent lower bound on the second bracket.
A strictly positive lower bound establishes an economic gain from the returned
policy even when its optimality certificate is conservative.

\section{The Trained Continuation as the Certified Bellman Object}
\label{sec:r41direct}
A learned continuation enters the economic argument only when the function
actually returned by training, rather than its training target, satisfies the
Bellman comparison. We now set $v_t=f_t$ in
Theorem~\ref{thm:r39comparison}, where $f_0,\ldots,f_T$ are the stored neural
networks. The continuation used inside the date-$t$ Bellman operator is
$f_{t+1}$. Its replacement by a reference approximation is not permitted in
this calculation. This construction completes the connection between the
learned scalar continuation and the returned dynamic policy; it does not
require the training procedure to have found a global minimizer.

\subsection{A certificate that does not use the training reference}
Let $\{x_i\}$ and $\{a_j\}$ cover $X$ and $A$ with radii $h_x$ and $h_a$,
respectively. Ties in nearest-node selection have a fixed measurable rule.
Suppose $f_t$ has Lipschitz constant $L_t$, and let $L_{x,t}$ and $L_{a,t}$
bound the variation of $Q_t(f_{t+1};x,a)$ in state and action. Verified
arithmetic returns
\[
 f^-_{ti}\leq f_t(x_i)\leq f^+_{ti},\qquad
 q^-_{tij}\leq Q_t(f_{t+1};x_i,a_j)\leq q^+_{tij}.
\]
These intervals enclose evaluations of the network coefficients and the
primitive transition and cost. Define $L_{ti}=\min_jq^-_{tij}$ and
$U_{ti}=\min_jq^+_{tij}$. Store any feasible grid action $a_{j(t,i)}$;
minimizing $q^+_{tij}$ is one implementable choice.

\begin{proposition}[Direct neural Bellman certificate]
\label{prop:r41direct}
Set
\begin{align}
 \ell_t&=\min_i(f^-_{ti}-U_{ti})-(L_t+L_{x,t})h_x,
       \label{eq:r41direct}\\
 u_t&=\max_i(f^+_{ti}-L_{ti}+L_{a,t}h_a)
                         +(L_t+L_{x,t})h_x,\nonumber\\
 \eta_t&=\max_i(q^+_{ti,j(t,i)}-L_{ti})+L_{a,t}h_a+2L_{x,t}h_x.
       \nonumber
\end{align}
If $\ell_T\leq f_T-g\leq u_T$ on $X$, the nearest-node actor $\pi$ satisfies
\begin{equation}
 0\leq J_0^\pi(x)-V_0^*(x)\leq
 E_f:=\sum_{t<T}B_t(u_t-\ell_t+\eta_t)+B_T(u_T-\ell_T).
 \label{eq:r41E}
\end{equation}
Its comparison class is all feasible adapted policies. If a separate
own-policy calculation gives $J_0^\pi(x)\leq\overline J_0^\pi(x)$, then
\begin{equation}
 J_0^\pi(x)-V_0^*(x)
 \leq\overline J_0^\pi(x)-f^-_0(x)+\sum_{t=0}^TB_tu_t.
 \label{eq:r41query}
\end{equation}
Neither certificate reads a reference spline, its nodal values, or its
Bellman brackets.
\end{proposition}
\begin{proof}
At a node the continuous-action minimum lies in
$[L_{ti}-L_{a,t}h_a,U_{ti}]$: every feasible minimizer is within $h_a$ of a
grid action. Subtraction gives the nodal residual interval. The difference
$f_t-T_tf_{t+1}$ has Lipschitz constant at most $L_t+L_{x,t}$, which proves
the first two bounds throughout each covering cell. Moving a chosen action
and the Bellman minimum from a node to the same state costs at most
$2L_{x,t}h_x$. This proves the actor bound. Apply
Theorem~\ref{thm:r39comparison}, whose lower value bound also proves
\eqref{eq:r41query}. The argument uses the stored neural functions at both
successive dates, including the terminal date.
\end{proof}

In contrast to \eqref{eq:r39mesh}, the state remainder includes $L_t$:
the trained network need not interpolate the verification nodes. Omitting
that term would incorrectly treat the training target as the learned
function. Additive shifts $f_t\mapsto f_t+k_t$ change both residual endpoints
by $k_t-\beta_tk_{t+1}$ and leave their oscillation and the actor allowance
unchanged. The terminal shift completes the telescoping sum. Thus both
\eqref{eq:r41E} and the exact version of \eqref{eq:r41query} are invariant
to the arbitrary levels of the learned continuations.

For the scalar economy, a stored ReLU network has rational breakpoints and
rational affine coefficients because its floating coefficients are exact
dyadic numbers when interpreted as stored data. On each activation region
we compute its affine coefficients from those weights. This is an exact
compilation of the trained network, not a fitted spline or a conventional
Bellman solution. If a breakpoint is rounded by at most $e_r$, choosing an
adjacent affine extension adds at most $2L_te_r$ to the value enclosure.
Rational coefficient enclosures and the interval width of the evaluation
argument are charged separately. Appendix~\ref{app:r41arithmetic} proves
this statement and specifies the arithmetic checks.

The terminal cost in \eqref{eq:r39capital} is piecewise quadratic. Hence the
extrema of $f_T-g$ occur at an activation breakpoint, at $3/8$, at a domain
endpoint, or at a stationary point in an intervening region. All these
points are calculated in rational arithmetic. Terminal fit error is thereby
included in the same neural chain, rather than set to zero because the
terminal primitive is known.

The deposited networks were trained on the earlier conventional nodal
values. That training-data cost remains part of their provenance. The present
verification changes the certified mathematical object, not that history:
there are no new neural training runs, and there is no assertion that the
reference was unnecessary for the original fit. The verification program
continues to run when the conventional spline evaluator is disabled.

\subsection{From the stored actor to its numerical execution}
\label{sec:r41deployment}
The stored actor is not the last object entering an economic calculation.
Let the observed state be rounded with error at most $s_x$. Suppose selected
actions and grid cutpoints are represented exactly, and their comparisons
are exact. Nearest-node selection at the rounded state can choose a different
cell, but its node is within $h_x+s_x$ of the true state. No continuity of
the tabulated actor is required. Suppose further that the numerical
transition and stage-cost graphs satisfy
\[
 \|\widetilde F_t(x,a,z)-F_t(x,a,z)\|\leq s_{F,t},\quad
 |\widetilde c_t(x,a)-c_t(x,a)|\leq s_{c,t},\quad
 |\widetilde g(x)-g(x)|\leq s_g.
\]
The bounds include conversion of their inputs and clipping. Both transitions
stay in $X$. The recursive risk functional itself retains its stated
mathematical meaning and is evaluated by enclosing intervals.

\begin{proposition}[Two execution contracts]
\label{prop:r41deployment}
Under Proposition~\ref{prop:r41direct}, the rounded-state actor
$\pi^{\rm exec}$ applied to the original economic transition satisfies
\begin{equation}
 0\leq J_0^{\pi^{\rm exec}}-V_0^*
       \leq E_f+\sum_{t<T}2B_tL_{x,t}s_x.
 \label{eq:r41actor_execution}
\end{equation}
For the numerical economy that also uses $\widetilde F_t$,
$\widetilde c_t$, and $\widetilde g$, its recursive cost instead satisfies
\begin{equation}
 \widetilde J_0^{\pi^{\rm exec}}-V_0^*
 \leq E_f+E_{\rm exec},\quad
 E_{\rm exec}:=\sum_{t<T}B_t
 (2L_{x,t}s_x+s_{c,t}+\beta_tL_{t+1}s_{F,t})+B_Ts_g.
 \label{eq:r41execution}
\end{equation}
The left side of \eqref{eq:r41execution} need not be nonnegative: its two
values belong to different transition and accounting models.
\end{proposition}
\begin{proof}
The nearest-node proof of Proposition~\ref{prop:r41direct} replaces $h_x$
by $h_x+s_x$ in the actor allowance. Feasibility is unchanged. This proves
\eqref{eq:r41actor_execution}. Monotonicity and cash invariance give
\[
 |\rho_t(f_{t+1}(\widetilde F_t))-
   \rho_t(f_{t+1}(F_t))|\leq L_{t+1}s_{F,t}.
\]
Thus the numerical chosen-action operator has an additional upper error
$2L_{x,t}s_x+s_{c,t}+\beta_tL_{t+1}s_{F,t}$ relative to the original Bellman
operator. Backward induction on the numerical policy, with terminal error
$s_g$, adds the discounted sum of these errors to the proof of
Theorem~\ref{thm:r39comparison}. The lower bound on the original optimum is
unchanged. This proves \eqref{eq:r41execution} without identifying a rounded
state trajectory with an exact one.
\end{proof}

We instantiate the second contract with separately rounded binary32
operations, round-to-nearest, and gradual underflow. On the specified domain
there is no overflow. With $u=2^{-24}$ and $\alpha=2^{-150}$, each elementary
rounding obeys $|\operatorname{fl}(y)-y|\leq u|y|+\alpha$. Rational
range-and-error propagation through the actual arithmetic graph gives
\[
 s_x<5.961\,10^{-8},\quad s_F<3.167\,10^{-7},\quad
 s_c<5.267\,10^{-7},\quad s_g<6.413\,10^{-7}.
\]
The stage-cost bound displayed here covers both prices. All action values
and cutpoints at the reported dyadic resolutions are exactly representable;
this accounts for zero action-storage error without setting the whole
execution allowance to zero. Clipping is nonexpansive. Fused evaluation,
flush-to-zero, altered state acquisition, and a physical actuator would
require their own contracts. The software experiment evaluates the rounded
transition and rounded costs, not a claim of physical deployment validation.

\section{Verified floating-point implementation and complete work}
\label{sec:r39floating}

The factor iteration is useful only if its whole-step allowance can be
obtained for the operations actually performed. The next result supplies such
an allowance without exact matrix products, exact solves, or a claim that a
small computed residual is automatically rigorous. The residuals themselves
must be enclosed.

\begin{proposition}[Residual verification of a native update]
\label{prop:r39floating}
Let $M$ be the fixed stored symmetric target with a proved bound $M\succeq mI$,
$m>0$, and let $W\in\R^{p\times d}$. Let $Z$ and $V$ be arbitrary stored
floating-point proposals. Define their exact-real residuals by
\[
 R=MZ-W^\top W,\qquad D=V-\tfrac12W(3I-Z).
\]
If outward arithmetic proves $\|R\|_F\leq r_R$, $\|D\|_F\leq r_D$, and
$\|W\|_2\leq w$, then the ideal update
$\overline W=\tfrac12W(3I-M^{-1}W^\top W)$ satisfies
\begin{equation}\label{eq:r39step}
 \|(V-\overline W)M^{-1/2}\|_2
 \leq \nu:=\{r_D+wr_R/(2m)\}/\sqrt m.
\end{equation}
Writing $E=M^{-1/2}W^\top WM^{-1/2}-I$, if $\|E\|_2\leq r<1$, then
\begin{equation}\label{eq:r39envelope}
 \|M^{-1/2}V^\top VM^{-1/2}-I\|_2
 \leq (3+r)r^2/4+2\nu+\nu^2.
\end{equation}
Thus $2\nu+\nu^2\leq(1-r)r^2/4$ certifies a squared-radius update.
Alternatively, a direct Gram enclosure can certify a terminal target without
certifying this contraction.
\end{proposition}
\begin{proof}
$Z-M^{-1}W^\top W=M^{-1}R$, so
$V-\overline W=D-\tfrac12WM^{-1}R$. Submultiplicativity gives
\eqref{eq:r39step}. Put $Y=WM^{-1/2}$ and $\overline Y=\overline WM^{-1/2}$.
Then $\overline Y=\tfrac12Y(3I-Y^\top Y)$ and
$\overline Y^\top\overline Y-I=(-3E^2+E^3)/4$.
For each eigenvalue $e\in[-r,r]\subset(-1,1)$, this polynomial is nonpositive,
and $(1+e)(1-e/2)^2\leq1$. Hence $\|\overline Y\|_2\leq1$.
The implemented perturbation adds at most $2\nu+\nu^2$ to the Gram error,
proving \eqref{eq:r39envelope}. The argument works for every $r<1$; it does
not apply the earlier $r\leq1/2$ budget outside its stated range.
\end{proof}

The implementation forms proposals in binary32 or binary64 and verifies them
with outward binary64 matrix balls. A cached approximate inverse is permitted
only for the identical fixed target and precision. Residual $R$ is evaluated
against the original stored target and factor, so conversion and solve errors,
including those of the cached inverse, enter the enclosure. A failed binary32
proposal can be retried in binary64. Exhausting the available native types is
reported as a failure, not as a proof that higher precision would fail.
Positive spectral bounds are certified by sufficient Gershgorin inequalities;
matrices failing this sufficient check are not silently accepted on the basis
of a computed eigenvalue. The experiments concern exact-real feedback rules
with stored binary coefficients. Floating-point execution of the actor needs
the separate execution allowance already present in the policy theorem.

\subsection{Work and storage accounting}
Let $n_b$ be the number of target-inverse constructions, $n_a$ the number of
proposed factor updates including rejected proposals, $n_e$ the number of
complete own-policy evaluations, $n_c$ the number of complete policy checks,
and $n_h$ the number of actor linear solves. For dense matrices, a sufficient
arithmetic-operation account for this implementation is
\begin{equation}\label{eq:r39work}
 O\!\left(n_bd^3+n_a(pd^2+d^3)
              +(n_e+n_c)Td^3+n_hd^3+\mathcal W_{\rm io}\right).
\end{equation}
Here $\mathcal W_{\rm io}$ includes recorded serialization and output work.
The ball products and residual checks have the same dense-operation order as
the proposed update, but their constants are not those of a single BLAS call.
A working storage bound, excluding the accumulated audit log, is
$O(T(pd+d^2)+pd+d^2)$ words. The stored factor remains binary64 even when its
proposal was calculated in binary32. Reduced proposal precision therefore does
not imply reduced persistent factor storage in this implementation.

For an accepted fixed-target contraction path, a sufficient update cap is
$K=\lceil\log_2(\log(1/\tau)/\log(1/r_0))\rceil_+$.
An ideal harmonic scalar initialization has
$r_0=(\kappa-1)/(\kappa+1)$, where $\kappa=L/m$; the implemented initial
radius includes rounding and is checked separately. At most two native
proposal attempts are made per accepted step when both types are available.
Conditioning also enters \eqref{eq:r39step} through $m^{-1}$ and $m^{-1/2}$.
The economic target determines $\tau$ through the previously derived actor,
evaluation and Gram allowances; the entropic moment margin enters those
allowances rather than disappearing into an iteration count. If a gate fails,
the realized counters, not $K$ alone, determine complete work.
Equation~\eqref{eq:r39work} is an arithmetic-operation bound at the stated
native precisions, not a total arbitrary-precision bit-complexity theorem.

For the scalar nonlinear construction, $N+1$ states, $A+1$ actions and $s$
shocks require $O(TNAs\log N)$ operations with binary-search spline evaluation,
in addition to the cost of certified transcendental evaluation. A sweep over
sorted successor states can change this bound but is not presumed here.
All fixed-rule shock-tree evaluations are also charged; their worst-case size
is $O(s^T)$. Neither this finite construction nor the endpoint verifier asserts
dimension-independent approximation complexity.

\subsection{Relation to numerical methods}
Classical polar and matrix-function iterations supply the algebraic update,
not its economic interpretation. Mixed-precision iterative refinement separates
working, correction and residual precision; \citet{CarsonHigham2018R39} provides
a systematic analysis of that separation. Verified numerical computation, as
surveyed by \citet{Rump2010R39}, makes the additional distinction between a
computed residual and an enclosing bound. Our residual construction uses both
ideas, while requiring neither to imply an economic policy guarantee by
itself. That guarantee comes from the Bellman comparison and its allocation
of evaluation, actor and continuation error. Likewise, the recursive-risk
quadratic comparator exploits the structural control solution, whereas
Section~\ref{sec:r39nonlinear} does not use a quadratic closure. The new
contribution is the connection between certified numerical work and a returned
policy in the original NBO framework, not a claim to have invented polar
iteration, spline approximation, or residual verification.

\section{Identified comparisons and a nonlinear economic counterfactual}
\label{sec:r39evidence}

The original exact-construction cases remain below, including the faster
structural comparator and the earlier adverse reuse comparisons. They are not
used to identify a precision effect. We instead analyze the deposited R38
factorial and nonlinear execution, and add an independently checked all-state
network certificate. No historical clock is replaced by a reconstruction.

\subsection{Design and accounting}
The fixed protocol contains 315 isolated services: 288 quadratic policy
services, 12 spline services, 12 trained-network services and three nonlinear
accuracy-frontier services. There are 105 deterministic economy--method objects
with three timing repetitions each. A separate warm-up is excluded from the
contrasts. Execution order is shuffled with a fixed seed, and the numerical
libraries use one thread. Neither CPU affinity nor frequency was controlled.
The displayed time is the complete child service through output synchronization;
the archive also records startup-inclusive parent time. Minima and maxima
accompany medians rather than a significance claim from three observations.
Memory is process high-water RSS, not exclusive algorithm memory.

The quadratic design crosses fixed binary64 versus adaptive native precision
with inverse rebuilding versus within-target caching. It adds a cached fixed
method selected from binary32 and binary64, in that order, and the structural
solution. All tuning trials, failed policy checks and precision retries enter
the service boundary. The selected fixed type is the lowest successful type in
that two-element set, not the minimum possible arbitrary mantissa length.
The update rule is held fixed in this factorial; comparisons with the older
update rules remain separately identified in the historical table.

Every candidate policy is checked initially and after an actor change; the
service stops at its first successful economic certificate. The target grid is
$10^{-2},10^{-4},10^{-6},10^{-8},10^{-10}$. Further objects vary $(d,T)$ over
$(8,4),(16,8),(32,4),(4,12)$ with condition parameters $8,16,16,4$, change the
valuation by factors $1.005,1/4,4$, include cold and transported starts, and
include an additional moment-margin case. This finite design tests several
operating regimes; it is not an asymptotic scaling experiment or an exhaustive
stress test near the entropic moment boundary.


% Materialized from revisions/2026-10-07-r41/manuscript/factorial_table.tex
\begin{table}[htbp]
\centering\small\setlength{\tabcolsep}{5pt}
\caption{Precision and caching at a common policy target}\label{tab:r39factorial}
\begin{tabular}{lrrrrr}
\toprule
Method & Median (s) & Range (s) & Updates & Builds & Loss bound\\
\midrule
Fixed64, rebuilt & 0.416 & [0.416, 0.420] & 11 & 11 & $5.59\times10^{-6}$\\
Fixed64, cached & 0.415 & [0.415, 0.421] & 11 & 3 & $5.59\times10^{-6}$\\
Adaptive, rebuilt & 0.421 & [0.418, 0.422] & 11 & 11 & $5.59\times10^{-6}$\\
Adaptive, cached & 0.433 & [0.417, 0.433] & 11 & 3 & $5.59\times10^{-6}$\\
Tuned native, cached & 0.424 & [0.423, 0.428] & 11 & 3 & $5.59\times10^{-6}$\\
Structural & 0.170 & [0.168, 0.173] & 0 & 0 & $4.46\times10^{-28}$\\
\bottomrule
\end{tabular}
\par\smallskip\begin{minipage}{.96\linewidth}\footnotesize $d=2$, $T=4$, target $10^{-4}$. Three isolated complete-service clocks per object. Builds count target inverses, not actor solves. All recorded failed checks and tuning are charged.\end{minipage}
\end{table}



The factorial removes the original inverse-caching confound. It does not show
a stable timing advantage from adaptive precision. At tolerance $10^{-4}$,
the fixed64 cached median is about $0.415$ seconds and the adaptive cached
median about $0.433$ seconds; the structural median is about $0.170$ seconds.
The two cached factor procedures each make eleven updates and three inverse
constructions. The uncached procedures construct eleven inverses. Reusing an
inverse changes this count, but verification and own-policy evaluation can
dominate the small-matrix service clock.


% Materialized from revisions/2026-10-07-r41/manuscript/accuracy_table.tex
\begin{table}[htbp]
\centering\small\setlength{\tabcolsep}{5pt}
\caption{First certified policy across accuracy targets}\label{tab:r39accuracy}
\begin{tabular}{rrrrrr}
\toprule
Target & Adaptive bound & Bound/target & Adaptive (s) & Fixed64 (s) & Structural (s)\\
\midrule
0.01 & 0.0023 & 0.23 & 0.352 & 0.353 & 0.175\\
0.0001 & $5.59\times10^{-6}$ & 0.0559 & 0.433 & 0.415 & 0.170\\
$1\times10^{-6}$ & $4.68\times10^{-11}$ & $4.68\times10^{-5}$ & 0.444 & 0.442 & 0.171\\
$1\times10^{-8}$ & $4.35\times10^{-11}$ & 0.00435 & 0.500 & 0.482 & 0.169\\
$1\times10^{-10}$ & $4.35\times10^{-11}$ & 0.435 & 0.486 & 0.491 & 0.175\\
\bottomrule
\end{tabular}
\par\smallskip\begin{minipage}{.96\linewidth}\footnotesize Factor methods use caching. Times are medians; complete dispersion and all methods appear in the supplement. Discrete oversolving is displayed, not suppressed.\end{minipage}
\end{table}



First-pass stopping removes unnecessary continuation after a successful policy
check. It does not force a discrete update to land at the target. The achieved
bound-to-target ratio for adaptive caching ranges from approximately
$4.68\times10^{-5}$ to $0.435$ over these five targets. In particular, the
$10^{-6}$ object still oversolves substantially. Table~\ref{tab:r39accuracy}
therefore reports achieved error as well as work; it is not presented as a
frontier with uniformly binding constraints. The structural solver also
oversolves, for a different reason: the economy permits an exact structural
policy. All 288 policy services certify their declared targets. Complete
counters, process memory and timing dispersion for every cell are supplied in
the supplement and the machine-readable audit.

Small and large transported changes have different operational consequences.
For adaptive caching, the $1.005$ change uses three new factor updates rather
than twelve from a cold start. The factor-$4$ change fails all three transported
factor gates and falls back to cold construction. The factor-$1/4$ change also
fails the transported gates. Setup of the old policy is charged jointly in the
reported warm service; accordingly its total clock need not beat a cold solve.
These are observed gate outcomes, not a success rate estimated for an unknown
population of future changes.

\subsection{Nonlinear approximation, policy accuracy and economic gains}
The economy in Section~\ref{sec:r39economy} is run for four decision dates at
$(P,\theta)\in\{1,4\}\times\{0,1\}$. The direct construction uses $N=256$
state cells and $A=128$ action cells. The trained continuation has 31 hidden
units, changes all affine hidden weights, and uses 300 Adam steps per date.
Training data are deterministic certified Bellman nodal values, not independent
observations. Every state and every continuous feasible action are covered by
the interpolation and action allowances in the theorem. The finite shock sum
is evaluated through interval operations with Taylor enclosures for exponential
and logarithm; no analytic Gaussian continuation is used.


% Materialized from revisions/2026-10-07-r41/manuscript/nonlinear_table.tex
\begin{table}[htbp]
\centering\small\setlength{\tabcolsep}{5pt}
\caption{Nonlinear capital: policy bounds and complete work}\label{tab:r39nonlinear}
\begin{tabular}{rrrrrr}
\toprule
Price & Risk & Neural bound & Spline bound & Neural (s) & Spline (s)\\
\midrule
1 & 0 & 0.2359 & 0.2352 & 2.846 & 0.240\\
1 & 1 & 0.2366 & 0.2359 & 5.011 & 1.311\\
4 & 0 & 0.2533 & 0.2520 & 2.999 & 0.409\\
4 & 1 & 0.2542 & 0.2528 & 6.222 & 2.564\\
\bottomrule
\end{tabular}
\par\smallskip\begin{minipage}{.96\linewidth}\footnotesize $T=4$, $N=256$, $A=128$. Cost includes reference construction and, for neural services, all fitting and verification. Price-four services also include construction and evaluation of the old rule. Bounds are in model cost units.\end{minipage}
\end{table}



The trained policies have global loss bounds between approximately $0.236$ and
$0.254$. These are deliberately reported in the model's cost units, not
reinterpreted as $10^{-4}$ certificates. The spline method is faster in all
four complete-service comparisons. A separate constructive accuracy sweep
has first-pass bounds $0.9352,0.4704,0.2359,0.1181$ at targets
$1,1/2,1/4,1/8$, respectively, using $N=64,128,256,512$ and $A=N/2$.
The corresponding bound-to-target ratios lie between $0.935$ and $0.946$.
Its cumulative construction clocks include the unsuccessful coarser meshes;
they are reported separately from the full service clock and are not four
independent experimental observations from each repetition.

The price change nevertheless supports a substantive policy comparison.
The old rule is the certified conventional rule constructed at $P=1$ and then
held fixed; the new rule is selected by the trained neural continuation at
$P=4$. Both are evaluated under the new price, by an independent finite-shock
tree enclosure. Table~\ref{tab:r39economic} reports the second bracket in
\eqref{eq:r39decomposition}. Its lower bound is positive at each reported
initial state for both risk specifications. Thus the new policy strictly
improves on retaining the old rule under the changed investment burden.
This comparison is not a claim that the neural rule outperforms the newly
optimized spline rule. The latter remains the stronger work comparator.


% Materialized from revisions/2026-10-07-r41/manuscript/economic_table.tex
\begin{table}[htbp]
\centering\small\setlength{\tabcolsep}{5pt}
\caption{Cost saved by reoptimization after the investment-price change}\label{tab:r39economic}
\begin{tabular}{rrrrrr}
\toprule
Risk & Capital & Old action & New action & Gain lower & Gain upper\\
\midrule
0 & 0.125 & 0.25000 & 0.21094 & 0.085570 & 0.085571\\
0 & 0.25 & 0.25000 & 0.16992 & 0.097610 & 0.099736\\
0 & 0.5 & 0.16797 & 0.10938 & 0.057079 & 0.057080\\
0 & 0.75 & 0.06836 & 0.05273 & 0.021237 & 0.030832\\
1 & 0.125 & 0.25000 & 0.21289 & 0.084553 & 0.084554\\
1 & 0.25 & 0.25000 & 0.17383 & 0.097370 & 0.098241\\
1 & 0.5 & 0.16797 & 0.10938 & 0.057067 & 0.057068\\
1 & 0.75 & 0.07031 & 0.05273 & 0.023443 & 0.028670\\
\bottomrule
\end{tabular}
\par\smallskip\begin{minipage}{.96\linewidth}\footnotesize Old: conventional rule constructed at price one, held fixed at price four. New: neural-selected rule at price four. Gains compare both rules under the new price. Displayed gain endpoints are rounded outward to six decimals; action entries are descriptive.\end{minipage}
\end{table}



Higher investment burden is accompanied here by lower initial investment in
low and intermediate capital states. The fixed-rule exposure formula explains
why keeping the old action rule incurs a positive repricing burden, while the
independent policy comparison quantifies how much of that burden is avoided.
The gain is economically interpretable even though the uniform global
optimality bound is more conservative. It does not identify an empirical
elasticity or a calibrated welfare effect.

\subsection{New all-state network verification}
The endpoint verifier in Proposition~\ref{prop:r39breakpoints} is applied to
all 60 deposited date-specific network records. Deduplication by coefficients
and reference function gives 17 distinct networks, not 60 independent fits.
Every new uniform network-to-spline bound is below $0.001$; the previous
mesh-Lipschitz checks were between $0.027$ and $0.031$. Independent exact
rational evaluation of every affine-piece endpoint verifies all 17 bounds,
with two further fixtures covering zero slopes and nondyadic roots.
This is an improvement in verification resolution, not a new timing experiment
or an automatic improvement in the previously recorded policy-loss numbers.
The audit records exact numerators and denominators, outward upper bounds,
source identities and the mapping from every service to its distinct network.

\paragraph{Interpretation.}
The results establish a constructive NBO comparison outside exact quadratics,
a residual-verified native factor implementation, and certified nonlinear
reoptimization gains. They also show why a certificate is not a work-superiority
result. The structural quadratic solver and the nonlinear spline baseline
remain favorable comparators. General diffusion, endogenous-preference,
temporal-self and game applications retain their original statements and
application-specific obligations; the two completed constructions are not
silently substituted for those distinct assumptions.

\section{Direct Neural Evidence and Activated Precision Adaptation}
\label{sec:r41evidence}
The direct certificate changes both the verified continuation and the
returned policy. At each date it evaluates the trained current network
against a Bellman update using the trained next network, then stores an
action minimizing an upper interval endpoint. The terminal fit is included.
We subsequently evaluate that newly constructed policy on its own shock
tree. Its displayed values are not values of the former spline policy or
of the previously selected neural actor.

\subsection{Resolution, policy accuracy, and the economic comparison}
The design retains the four investment economies, with prices $P=1,4$ and
risk parameters $\theta=0,1$. It uses the frozen trained networks and tries
state--action resolutions $(256,128)$, $(1024,512)$, $(2048,1024)$, and,
if necessary, $(4096,2048)$. Refinement stops at the first mesh at which the
maximum certified loss at the four displayed initial states is at most
$0.02$. The first two meshes fail that query target in every economy;
the third passes in every economy. These are twelve new post-training
verification services, not twelve new training runs.

\begin{table}[htbp]\centering\small
\caption{Direct certificates for the stored neural continuations}\label{tab:r41direct}
\begin{tabular}{rrrrrr}\toprule
$P$ & $\theta$ & All-state bound & Query bound & Execution allowance & All refinements (s)\\\midrule
1 & 0 & 0.043092 & 0.016279 & 9.719 & 8.978\\
1 & 1 & 0.043334 & 0.016402 & 9.751 & 32.091\\
4 & 0 & 0.045472 & 0.017829 & 10.235 & 8.464\\
4 & 1 & 0.045813 & 0.017781 & 10.266 & 31.319\\
\bottomrule\end{tabular}
\par\smallskip\begin{minipage}{.97\textwidth}\footnotesize Notes: The query column is the maximum certified loss at $x=1/8,1/4,1/2,3/4$, not an all-state bound. The execution allowance is in units of $10^{-6}$ and covers rounded state acquisition, transition, and costs. Every final mesh has 2,048 state cells and 1,024 action cells; the clocks sum all three attempted meshes through the result-file flush and synchronization. They exclude the inherited training service and are not an end-to-end speed comparison.\end{minipage}\end{table}

The all-state loss remains between approximately $0.0431$ and $0.0459$.
The sharper displayed-state bounds, below $0.0179$, use independently
computed own-policy upper values and the direct neural lower bound on the
optimal value. They must not be extrapolated to unreported states. The
nonzero numerical-execution allowance is approximately $10^{-5}$; it is
added to the appropriate global comparison, not silently discarded.

At the accepted resolution we independently construct a conventional
spline policy and evaluate it at the same initial states. The supplement
reports each interval for neural cost minus this newly optimized spline
cost. All sixteen lie inside $[-0.001,0.001]$. Fourteen have strictly positive
lower endpoints; none has a strictly negative upper endpoint, and two
contain zero. Thus a small descriptive cost margin is established at these
states, while the conventional policy still has strictly smaller verified
cost in fourteen comparisons. Interval overlap is not used as a test of
equivalence. The margin is an explicit descriptive accounting tolerance,
not a preregistered economic indifference threshold or a population claim.

For $P=4$, gains over continuing the former $P=1$ rule have positive lower
endpoints at all eight displayed state--risk pairs, ranging from about
$0.0233$ to $0.0988$. The largest displayed-state near-optimality upper
bound is smaller than the smallest of these gain lower bounds. This
establishes that the direct neural policy's certified regret at the reported
states is smaller than the verified benefit of reoptimization there. It
does not make the price change, the benefit of reoptimization, or the
conventional competitor specific to NBO. All earlier adverse structural,
spline, and finite-law comparisons remain in the article and companion.

\subsection{When binary32 is no longer sufficient}
The earlier factorial correctly separates precision from caching, but its
adaptive paths never require binary64. A separate pressure catalogue now
crosses dimensions $2,8$, target condition numbers $1,16$, and policy targets
$10^{-10},10^{-12},10^{-14}$. The horizon is four, innovation scale is
$0.001$, and risk sensitivity is $0.1$. The inherited acceptance gates are
unchanged. Every proposed factor records its precision, input and output
hashes, acceptance decision, and target identity. A rejected binary32
proposal is charged even when a binary64 proposal subsequently succeeds.

\begin{table}[htbp]\centering\small
\caption{Precision pressure: accepted updates and rejected attempts}\label{tab:r41precision}
\begin{tabular}{lrrrrr}\toprule
Method & Certified / 12 & Accept 32 & Accept 64 & Reject 32 & Reject 64\\\midrule
adaptive cached & 12 & 116 & 42 & 42 & 0\\
fixed32 cached & 2 & 49 & 0 & 10 & 0\\
fixed64 cached & 12 & 0 & 158 & 0 & 0\\
tuned cached & 12 & 49 & 147 & 10 & 0\\
\bottomrule\end{tabular}
\par\smallskip\begin{minipage}{.97\textwidth}\footnotesize Notes: Twelve objects cross dimension and conditioning with three economic targets, $10^{-10}$, $10^{-12}$, and $10^{-14}$. Failed fixed32 services are retained. Counts for tuned precision include unsuccessful trials before restarting at the higher precision. One execution per object identifies operation of the mechanism, not a timing distribution.\end{minipage}\end{table}

Ten of the twelve adaptive services use both precisions; the other two
succeed using binary32 alone. Fixed binary32 certifies only two objects.
Adaptive and fixed binary64 each certify all twelve. The adaptive path
accepts 116 binary32 and 42 binary64 proposals and rejects 42 binary32
proposals. These observations identify actual precision escalation rather
than inferring adaptation from the controller's name. They do not establish
that adaptive complete work is smaller than the best fixed native precision:
extra attempts and verification have real costs, and this pressure catalogue
has one clock per object. A stable efficiency ranking requires a separate
matched-work study.

\subsection{Provenance and what the work comparison measures}
The source-and-result records identify the R39 reviewed article and the
partially executed R40 development attempt separately. The latter stopped
when a deserialized policy's list-valued knots were used as a numerical
array. The new driver restores both knots and actors to arrays before
repricing the old policy. It reruns the entire catalogue in a fresh directory;
it neither overwrites the failed-run outputs nor substitutes new clocks
for the 315 earlier services. The inherited protocol predates this replay,
but some R40 outcomes had already been observed. We therefore do not call
the replay a new blinded or independently prospective experiment.

All three attempted mesh clocks, failed precision attempts, interval
verification, independent queries, result serialization, and the result
file's synchronization are retained. CPU affinity and one-thread execution
are recorded; frequency is not controlled. The inherited reference-label
construction and neural training costs are not part of the new
post-training clocks and cannot be forgotten in a from-scratch comparison.
The conventional comparison is made at the accepted mesh, not at a
separately optimized common stopping target. Accordingly neither set of
new clocks is a claim of neural end-to-end acceleration.

Theorem~\ref{thm:r41dimension} in the supplement gives the explicit cost of
multidimensional verification and continuous-innovation quantization,
including the additional own-policy query tree. This theorem extends the
certificate's domain and makes its dimensional cost visible. The new
executed nonlinear example remains one-dimensional. The evidence therefore
establishes a direct neural continuation certificate and an activated native
precision mechanism, not yet a competitive high-dimensional nonlinear
performance frontier.

\section{Retained construction evidence and earlier economic comparisons}\label{sec:r37evidence}
This section retains the earlier exact-construction evidence and adverse comparisons. The eight-service timing contrast changes precision and caching jointly and is descriptive, not a causal precision comparison. The identified current comparisons are in Section~\ref{sec:r39evidence}.

The evidence distinguishes local decisions, full policies, and the work of
constructing their continuations. The complete previous article and its
technical supplement contain the original tables, including unsuccessful
checks and comparisons unfavorable to NBO. The applications companion
retains every economic application. We first locate the full-policy results
relative to that evidence and then examine the precision-adaptive construction.

\subsection{From current decisions to full policies}
The original changed-future experiment contains 168 complete services and
840 future-specific outcomes. It evaluates current scalar actions conditional
on specified future policies. Unchanged reuse certifies 72 of its 120 outcomes;
check-and-refresh and fresh neural refitting each certify all 120. The 48
failed reuse checks are retained. Adaptive neural reuse costs more than fresh
neural refitting in every one of the eight dimension--volume cells, with
ratios between approximately 1.39 and 1.59. Cached sample-average optimization
is the least-cost fully certified service in all eight cells. The signed
withdrawal responses to stronger future production and valuation are valid
finite-law economic results, but full-support enumeration obtains the same
target more sharply and cheaply. These results support the decision
certificates; they are not evidence that the neural representation is
necessary for those counterfactuals.

The subsequent full-policy construction changes the comparison target, not
the meaning of those earlier observations. It uses vector controls,
continuous Gaussian innovations, and own-policy evaluation over 24 decision
dates, with matrix inequalities covering all real states. The additive
catalogue records 480 certified services and 608 failed intermediate checks
among 1,088 attempts. The recursive-risk catalogue records 72 certified
services under its retained execution identity. The latter catalogue spans
two state dimensions, three risk settings, three economic regimes, two
registered seeds, and two procedures. Its total recorded service work is
25.4238 seconds, with closing-ledger work reported separately. These are
finite-design records, not estimates of a success probability over an
unspecified population of training runs. Earlier recursive experiments
that failed a domain or precision check remain preserved; later successful
certificates do not erase those failures.

The full-policy comparison has no sampled-state coverage term: the coefficient
inequalities hold on all states. Its policy-class approximation term is zero
because the comparator is all adapted policies of finite cost. Its transfer
term is zero for the directly specified discrete economy, not for an
unverified approximation of the original diffusion. Evaluation, training,
actor, and action-execution allowances remain distinct. A successfully
verified capital policy therefore instantiates a full-policy result without
asserting a certificate for a separate game or utility specification.

\subsection{A frozen comparison of continuation constructions}
The new comparison holds the economic primitives and policy tolerance fixed
and varies only the construction procedure. It has two state and action
coordinates, four decision dates, Gaussian covariance $I/1000$, discount
factor $0.95$, recursive risk parameter $0.1$, and identity action and terminal
cost matrices. The transition matrix is
\[
 A=\begin{pmatrix}0.1&0.02\\0&0.11\end{pmatrix},\qquad B=I.
\]
The anchor state cost is $Q=I$; the changed-valuation state cost is
$\operatorname{diag}(201/200,199/200)$. Initial states satisfy
$\|x_0\|^2\leq2$, and the requested full-policy loss is $10^{-4}$.
Each method solves its own future backward. In the changed economy each
neural procedure starts from its own actual anchor factors, not from a
competitor's solution or an optimal-policy label.

The neural comparisons retain rectangular three-by-two factors,
noncommuting initial Gram matrices, nonzero evaluation radii, and perturbed
actor solves. Linear warm updates and fixed-precision quadratic refresh use
40 fractional bits. Precision-adaptive refresh uses the first integer
satisfying \eqref{eq:r37bits}, subject to the registered caps of 128 bits and
32 updates per fit, and builds the inverse of each fixed target at most once.
The structural comparator directly solves the same quadratic economy without
factor training. All methods receive the same analytic Gaussian information.
The terminal transformed continuation is evaluated directly for every method.

The complete source and protocol were committed before this eight-service
execution. Method order, regime order, precision limits, stopping conditions,
and failure handling were fixed. A failed service would have been written
and the execution terminated; no partial execution would have been reported
as eight successful services. Every returned policy is checked independently
by exact Bellman matrix subsolutions, in addition to its training and actor
acceptance conditions. The state and action coverage is algebraic, not an
interpolation claim inferred from the two-dimensional examples.


% Materialized from revisions/2026-10-07-r37/results/construction_table.tex
\begin{table}[t]
\centering
\caption{Complete construction at a common policy tolerance}
\label{tab:r37construction}
\small\setlength{\tabcolsep}{3.2pt}
\begin{tabular}{llrrrrr}
\toprule
Construction & Economy & Updates & Inverses & Bits & Gap bound & Seconds\\
\midrule
Linear warm & Anchor & 102 & 0 & 4080 & $1.83\!\times\!10^{-14}$ & 0.2808\\
Linear warm & Changed & 66 & 0 & 2640 & $1.88\!\times\!10^{-14}$ & 0.2359\\
Quadratic, fixed & Anchor & 6 & 6 & 240 & $1.53\!\times\!10^{-12}$ & 0.2224\\
Quadratic, fixed & Changed & 3 & 3 & 120 & $1.01\!\times\!10^{-9}$ & 0.2250\\
Quadratic, adaptive & Anchor & 6 & 3 & 102 & $1.89\!\times\!10^{-12}$ & 0.2006\\
Quadratic, adaptive & Changed & 3 & 3 & 54 & $9.90\!\times\!10^{-10}$ & 0.1971\\
Structural & Anchor & 0 & 0 & 0 & $0$ & 0.0396\\
Structural & Changed & 0 & 0 & 0 & $0$ & 0.0096\\
\bottomrule
\end{tabular}
\par\smallskip\begin{minipage}{0.98\linewidth}\footnotesize
\emph{Notes:} All eight policies satisfy the tolerance $10^{-4}$ uniformly on $\|x_0\|^2\leq2$. Inverses count constructions for training targets only. Bits is the sum of fractional precisions across hidden updates, not total arithmetic work. Complete seconds include evaluation, training, checks, actor solves, serialization, and durable detailed output. Common startup and the final ledger are unallocated. These are deterministic construction objects, not population samples.
\end{minipage}
\end{table}



All eight services satisfy the common policy tolerance. The linear procedure
uses 102 and 66 hidden updates in the anchor and changed economies. Both
quadratic procedures use six and three. The adaptive procedure reduces the
sum of stored fractional precisions from 240 to 102 in the anchor and from
120 to 54 in the changed economy. Its largest precision is respectively
21 and 18 bits. For the anchor it constructs three training target inverses,
rather than six; for the changed economy both quadratic procedures construct
three because each trained date needs only one update. The target-inverse
applications are still performed at every update.

In these two specified services the adaptive construction also has lower
recorded complete time than the two inherited neural procedures. The
structural procedure remains faster than every neural procedure and returns
an exact zero policy-gap bound. This outcome is informative: precision and
reuse improve the implemented neural construction, while the economy's
structure still makes direct control particularly inexpensive. We do not
convert a comparison between three neural implementations into superiority
over all candidate generators, or pool these deterministic objects with
historical fitting seeds to form a population test.

The service clock begins before primitive allocation and ends after the
detailed record has been serialized and durably written. It includes
own-policy evaluation, gates, all hidden updates, actor solves, exact policy
verification, and output. Common interpreter startup and the final summary
ledger are separately identified and not allocated selectively to a method.
The reported inverse counts concern training targets only; Gaussian evaluation
and actor solves are separately counted and included in the service clock.
The fractional-precision sum is not a count of all rational bit operations.
Process high-water memory is recorded as a process quantity, not as a
method-specific memory frontier. No simulation-transition advantage is claimed:
all eight services use analytic Gaussian evaluation and have zero simulated
transitions.

The comparison supplies an executed instance of
Corollary~\ref{cor:r37policy}: a trainable continuation, constructed with
explicit finite precision, yields a policy certified against full adapted
controls in its stated economy. The deterministic theorem applies to every
input meeting its basin, moment, and error conditions. The measured speed
comparisons apply to the eight executed services. Their different scopes are
part of the numerical conclusion rather than qualifications hidden from it.

% R15 component; only input paths and enumerated location/grammar prose are edited.
% Source: revisions/2026-10-04-r15/manuscript/economic_scope.tex; commit 1cb3cc9efe135966ec228dfaf51c6841a6dfc97c.
% Generated by code/integrate.py; edit extraction rules there.
% Reviewed source: revisions/2026-10-04-r14/manuscript/economic_scope.tex, line 1, commit f5021cefa71492babfcbfa580e0c984f59a9de26
\section{Economic Applications of the Bellman Structure}\label{sec:economic-scope}
The motivating economic problem combines external choices with the costly adjustment of preferences. Recursive utility changes how a given policy is evaluated; preference adjustment adds an internal choice to its Hamiltonian; and strategic interaction changes whose deviations define improvement. These distinctions determine the economic interpretation of NBO. The original diffusion applications and the directly specified recursive capital economy have distinct verification accounts. The latter provides the constructive full-policy instance studied above; it is not an unverified discretization certificate for the former. The preference, temporal-self, and game applications establish the corresponding economic objects and their own verification results.

\subsection{Costly preference adjustment and portfolio choice}
Let $u$ denote risk aversion and $X$ financial wealth. The household chooses consumption $c$, portfolio share $p$, and preference-adjustment effort $\theta$. The preference state is reflected on a declared interval $[\underline u,\bar u]\subset(1,\infty)$, and wealth follows the portfolio diffusion:
\begin{align}
 du&=\theta\,dt+\sigma_u\,dW^u+dL^{\underline u}-dL^{\bar u},\nonumber\\
 dX&=\{[r+p(\mu-r)]X-c\}\,dt+p\sigma_X X\,dW^X,
 \qquad d\langle W^u,W^X\rangle_t=\varrho\,dt.\label{eq:scope-preferences}
\end{align}
For a fixed consumption unit $\bar c>0$, flow utility is
\[
 U(c,u)-\tfrac12k\theta^2,
 \qquad U(c,u)=\frac{(c/\bar c)^{1-u}}{1-u}.
\]
The objective discounts both felicity and adjustment costs at rate $\rho$ and includes the specified wealth-exit and terminal payoff. The baseline uses $\bar c=1$. Because the household controls $u$, the cardinal normalization across preference states is a primitive: adding a function of $u$ changes the incentives to adjust preferences. The reflected state, positive consumption domain, and stopping data give the evaluation equation its economic and mathematical domain.

At an interior optimum with $V_X>0$ and $V_{XX}<0$, the three policy conditions are
\begin{align}
 c^*&=\bar c(\bar c V_X)^{-1/u}, &
 \theta^*&=V_u/k,\nonumber\\
 p^*&=-\frac{V_X(\mu-r)}{XV_{XX}\sigma_X^2}
       -\frac{\varrho\sigma_u V_{uX}}{XV_{XX}\sigma_X}.\label{eq:scope-preference-policies}
\end{align}
Box constraints require their corresponding feasible maximizers. The cross derivative in the portfolio rule links the external risk choice to the internal preference state. Its sign is an equilibrium object to be determined by evaluation. The adjustment rule likewise relates effort to the shadow value of changing preferences. A higher adjustment cost lowers the value of every fixed policy and hence the optimized value; whenever an envelope derivative exists,
\begin{equation}\label{eq:scope-preference-envelope}
 \partial_k V^k(t,u,X)
 =-\tfrac12\E^{a^k}_{t,u,X}
   \int_t^\tau e^{-\rho(s-t)}\theta_s^2\,ds\leq0.
\end{equation}
This comparative static holds with the feasible class and stopping payoff fixed.

The numerical application implements these choices jointly. The finite-action study evaluates complete policies by independent backward recursion. The continuous-action study trains the actor through its own payoff and verifies its hybrid policy with a complete action-box enclosure. Fresh training on nested grids separates the effect of resolution from the evaluation of a frozen coarse policy. A model-specific compensating-consumption calculation translates its cardinal utility bounds into an explicit transfer experiment. Application Sections~\ref{sec:ndu}, \ref{sec:ndu-neural}, and~\ref{sec:r9results} give the full model, all recorded policies, refinement evidence, and compensation account.

\subsection{Recursive utility and temporal selves}
Recursive utility enters policy evaluation through the policy's own continuation utility. With $a=1-1/\psi$, $b=1-\gamma$, and $bV>0$, the Epstein--Zin aggregator used here is
\begin{equation}\label{eq:scope-recursive}
 f(c,V)=\frac\rho a\{c^a(bV)^{1-a/b}-bV\}.
\end{equation}
The invariant utility interval and boundary payoff are part of the evaluation problem. Application Section~\ref{sec:recursive} derives the homothetic policy conditions and verifies the matching-bequest experiment, including the barriers that keep every policy utility in the aggregator's domain. This application tests recursive evaluation while maintaining the same separation between value approximation and feasible action improvement.

Time inconsistency requires a different improvement criterion. For the continuous discount kernel
\[
 h(s)=\beta e^{-\rho s}+(1-\beta)e^{-(\rho+\kappa)s},
 \qquad 0<\beta\leq1,\quad\kappa>0,
\]
two critics evaluate the same future policy at the two discount rates. Writing these values as $W_1,W_2$, the current self uses $V=\beta W_1+(1-\beta)W_2$ and chooses an instantaneous deviation by maximizing $u(a)+\LL^a V$. Application Section~\ref{sec:temporal} states the complete component equations and derives the spike-deviation condition. In the homothetic consumption economy, the equilibrium consumption share is the unique root in its admissible interval, and stronger present bias raises that share. Independently optimized finite-duration deviations check the limiting equilibrium condition. Thus the actor's objective follows from the temporal equilibrium concept, while the critics retain their policy-evaluation role.

\subsection{Capital investment and strategic deviations}
In the dynamic Cournot application, firm $i$ chooses investment in its own capital, with
\[
 dK^i=(I^i-\delta K^i)\,dt+\sigma_KK^i\,dW^i,
 \qquad q_i=AK^i.
\]
Its flow profit is $P_N(K)q_i-I^i-b_i(I^i)^2/2$, where the market-size normalization in $P_N$ is specified for each economy. Given rivals' policies, a critic evaluates the resulting own-firm payoff and the actor changes only its own action. Consequently, the relevant economic error is the gain from a complete unilateral deviation against the frozen rival profile.

Application Sections~\ref{sec:games} and~\ref{sec:dynamic-computation} specify the boundaries, investment sets, and market normalizations, and report independent full dynamic best responses to every principal finite neural profile. The resulting profit-unit bounds include all dates and states of the stated finite game. They also cover randomized and history-dependent unilateral deviations against those fixed Markov rivals, by finite dynamic programming. The record reports the cost of the exhaustive training oracle and any nonzero one-step equilibrium defects along with the final exploitability calculation.

These applications give the common method a precise economic interpretation. A recursive critic evaluates a utility process; preference and capital actors allocate resources and adjustment effort; a temporal actor satisfies a spike condition; and a strategic actor is assessed against unilateral best responses. Each application supplies the boundary conditions, approximation class, and error account required by its own economic question.

\section{Conclusion}\label{sec:r37conclusion}
Neural Bellman Operators separate policy evaluation from feasible improvement
while preserving the economic meaning of the continuation being learned.
Centered action errors identify what matters for a current decision;
own-policy Bellman comparisons identify what is required for a dynamic
policy. The constructive results connect these objects through updated
hidden weights, primitive error allocations, and explicit finite arithmetic.
They permit changing futures, noncommuting warm starts, two-sided
continuation errors, and inexact actor equations within the stated budget.

Precision-adaptive refresh gives a finite learning schedule whose relative
Gram envelope contracts quadratically and whose stored fractional precision
is allocated to the economic tolerance. A fixed target inverse is reused
within a fit, but evaluation of a new future, inverse applications, actor
solves, and independent policy checks remain part of the computation. The
new execution demonstrates this construction against full adapted-policy
comparators rather than only a current action catalogue. It also shows why
construction counts and complete economic work must be reported separately:
the earlier adaptive-cached implementation has smaller recorded clocks in its two exact regimes, whereas the structural procedure is faster. The factorial now separates those mechanisms and does not show a stable adaptive-precision advantage.

The controlled-diffusion formulation and the applications to recursive
utility, endogenous preferences, temporal selves, and games remain part of
the same research program, with their respective domains and comparison
conditions. The recursively risk-sensitive capital economy supplies a
complete constructive instance; it does not substitute its assumptions for
those of every other application. The resulting contribution is a proved
and implemented path from neural continuation construction to economic
policy accuracy, accompanied by transparent evidence about when the
particular construction does and does not improve measured work.


The nonlinear construction establishes the same decision-to-policy logic without quadratic closure. Its endpoint verifier certifies trained continuations on every state, and its investment experiment separates mechanical repricing from independently verified gains due to reoptimization. The structural and spline comparisons remain visible. These results extend the NBO paper through a second constructive economic realization while keeping numerical certification, learned-policy acceptance and empirical work superiority logically distinct.

The direct neural chain now certifies the learned continuation itself,
including its own terminal discrepancy and the policy selected from its own
future values. Its execution allowance is nonzero, and its pressure paths
actually reject binary32 before accepting binary64. These are constructive
advances in the same NBO framework. The direct comparisons also identify
what they do not establish: cost closeness at specified states is not neural
superiority, and an explicit multidimensional work theorem is not an
executed high-dimensional frontier.

\bibliographystyle{ecta-fullname}
\bibliography{revisions/2026-10-07-r41/ECTA_references}
\end{document}
